\section{Temporal Folding and Geometry}
\label{app:folding}

\subsection{Temporal-folding definition and geometry}
\label{app:folding-definition}

For an input $X\in\mathbb{R}^{C\times T}$ and a divisor $P$ of $T$, the
phase-interleaved adapter produces $F_P(X)\in\mathbb{R}^{CP\times T/P}$
according to
\[
F_P(X)_{cP+p,w}=X_{c,wP+p},\qquad 0\le p<P.
\]
This is a bijective permutation of the $CT$ samples, so folding changes only
the representation geometry. In contrast, a contiguous reshape assigns each
row a consecutive time chunk; both layouts have identical dimensions but
different local neighborhoods. The adapter pads only the bottom and right
edges to satisfy the backbone stride. For stride $s$, the padded dimensions
and non-padding height occupancy are
\[
H_{\rm pad}=s\left\lceil\frac{CP}{s}\right\rceil,\qquad
W_{\rm pad}=s\left\lceil\frac{T/P}{s}\right\rceil,\qquad
\rho_H=\frac{CP}{H_{\rm pad}}.
\]
Increasing $P$ raises vertical occupancy but reduces horizontal temporal
resolution and increases phase-wrap distance. The geometry rule therefore
selects the smallest valid $P$ reaching the target folded height; the complete
fold-factor sweeps are reported in Appendix~\ref{app:p-ablation}.

\paragraph{Exact invertibility.}
The mapping is bijective: for every original sample index
$t=wP+p$, the inverse is
$X_{c,t}=F_P(X)_{cP+(t\bmod P),\lfloor t/P\rfloor}$. Consequently,
$\|F_P(X)\|_0=\|X\|_0$ and no interpolation, averaging, or frequency-domain
operation is introduced. Any difference between temporal folding and contiguous
reshape therefore comes solely from the neighborhoods presented to the
backbone. The examples below visualize this change of coordinates; they are
not claims that the highlighted patterns are uniquely physiological or
causal.

\subsection{Visual interpretation and examples}
\label{app:folding-visuals}

The visualizations provide three complementary checks. First, the waveform
panels establish the original channel--time coordinates. Second, the folding
panels show that phase $p$ becomes a vertical coordinate while $w$ advances
by $P$ samples horizontally. Third, the real and synthetic examples verify
that localized transients, oscillations, and repeated motifs remain present
after the permutation. Pixel values should therefore be compared through the
index relation above, rather than by interpreting the folded image as a new
measurement.

We use the figures in a controlled-to-real progression. The simulated signals
in Fig.~\ref{fig:synthetic-folding-patterns} are shown from side and top views
before their phase--time rendering; this makes the two-dimensional motifs a
coordinate change rather than a new signal. Figure~\ref{fig:seizure-folding}
then shows the same correspondence for a real ictal CHB-MIT segment.

\begin{figure*}[t]
\centering
\includegraphics[width=0.94\textwidth]{side_top_view_folding.pdf}
\caption{The same EEG can be read from two complementary viewpoints. In the
side view (left), a human follows amplitude as a continuous function of time.
After lossless $P=4$ folding (right), the identical samples form a
two-dimensional phase--time texture: samples with phases $p=0,1,2,3$ become
adjacent vertical rows, while each column advances by four original samples.
The upper row is a controlled spike--slow-wave example and the lower row is a
real five-second F7--T7 segment from a CHB-MIT seizure recording. The arrows
indicate a change of representation and local adjacency, not a change in the
signal itself.}
\label{fig:side-top-view}
\end{figure*}
\FloatBarrier

\begin{figure*}[t]
\centering
\includegraphics[width=0.94\textwidth]{synthetic_folding_patterns.pdf}
\caption{Controlled simulated morphologies under lossless $P=4$ folding.
Each row presents the same waveform as a side view, a top view, and a
phase--time map. Spike--slow-wave complexes, recruited rhythmic oscillations,
and periodic discharges become localized two-dimensional motifs while every
sample is retained. These simulations explain the geometry; they are not
diagnostic templates.}
\label{fig:synthetic-folding-patterns}
\end{figure*}
\FloatBarrier
\begin{figure*}[t]
\centering
\includegraphics[width=0.94\textwidth]{seizure_lossless_folding.pdf}
\caption{A real seizure example processed by the visual interface. The
10-second CHB-MIT segment is labeled ictal and is resampled from 2560 to 2000
samples before clipping and scaling. Panel A shows all 16 bipolar channels;
panel B enlarges F7--T7 in the rhythmic interval; panel C shows the exact
$64\times500$ tensor after $P=4$ folding; and panel D expands its four phase
rows. All 32,000 values are retained.}
\label{fig:seizure-folding}
\end{figure*}
\FloatBarrier

\subsection{What a convolution sees}
\label{app:convolution}
Within one electrode's phase block, a vertical step changes the phase index,
while a horizontal step advances by $P$ original samples. A $k\times k$
convolution therefore mixes offsets
\[
\Delta t=P\Delta w+\Delta p,
\]
where $(\Delta w,\Delta p)$ is a kernel offset in folded coordinates. The
local patch is flattened and multiplied by the kernel weights, followed by
bias and nonlinearity; no special EEG operation is introduced. For $P=4$ and
$k=3$, the direct temporal offsets are
$\{-5,-4,-3,-1,0,1,3,4,5\}$.

At phase-block boundaries a vertical neighborhood can also cross electrodes,
so the interpretation depends on channel ordering. Figure~\ref{fig:channel-crossing}
makes the two regimes explicit. Inside one block, the three columns of a
$3\times3$ window are $P$ samples apart, so the window samples the waveform
like a dilated (atrous) pattern with rate $P$: for $P=4$ its nine taps cover
three consecutive samples around each of $t-P$, $t$, and $t+P$, which is
exactly the offset set above. At a block boundary, however, the $\Delta p=+1$
row is no longer the next phase of the same electrode: row $cP+P-1$ is
adjacent to row $(c+1)P$, so that row reads electrode $c+1$ instead, and its
offsets are measured in the neighbour's own time frame. The crossing therefore
mixes two electrodes only at block boundaries, and the induced neighborhood
depends on channel ordering.

\begin{figure*}[t]
\centering
\includegraphics[width=0.94\textwidth]{channel_crossing.pdf}
\caption{A $3\times3$ window at a phase-block boundary ($P=4$, real TUEV
segment; channels T5--O1 and FP2--F8). (a) Each channel occupies $P$
consecutive phase rows, separated by the bold line. The solid window straddles
the boundary, so its $\Delta p=+1$ row belongs to the next electrode, whereas
the dashed window for comparison stays inside a single electrode. (b) The same
boundary window on the original traces: one row of the window reads $c+1$
(blue), and within each electrode the taps form a dilated pattern with rate
$P=4$, matching $\Delta t=P\Delta w+\Delta p$.}
\label{fig:channel-crossing}
\end{figure*}
\FloatBarrier



\section{Detailed Results}
\label{app:detailed-results}

This appendix follows the per-dataset comparison style of the CBraMod and
REVE benchmark tables. All local runs inherit the fixed split code of the
applicable CBraMod or REVE benchmark and use the same subject- or trial-level
train, validation, and test partitions as the corresponding reference rows;
the remaining published rows are transcribed from these benchmark tables as
contextual references.

\paragraph{Channel ordering and montage preprocessing.}
All local vision models receive channels in the order produced by the
dataset-specific preprocessing pipeline. We do not reorder channels at runtime
by 10--20 position, Euclidean electrode distance, or any other
physical-coordinate rule, and the folding adapter does not insert or remove
channels. Consequently, the row-block order in Eq.~\ref{eq:fold} is exactly
the stored array order for each dataset. This convention makes the induced
cross-channel neighborhoods reproducible, although their physical
interpretation remains dataset-dependent. The channel lists and preprocessing
choices are included in the released preprocessing configuration.
Binary tasks report BA / PR-AUC / ROC-AUC;
multiclass tasks report BA / Cohen's $\kappa$ / weighted F1. Published values
are transcribed from the CBraMod and REVE tables, whereas local vision rows
are validation-selected results with one final test evaluation per seed.
EfficientNet-B0 uses its ImageNet initialization; ConvNeXt-Tiny and ViT-Small
use DINOv3 pretrained weights.
Supervised EEG models, EEG foundation models, and vision models are separated
into groups; the supervised rows reference EEGNet \citep{lawhern2018eegnet},
EEGConformer \citep{song2023eeg}, SPaRCNet \citep{jing2023development},
ContraWR \citep{yang2021self}, CNN-Transformer \citep{peh2022transformer},
FFCL \citep{li2022motor}, and ST-Transformer \citep{song2021transformer}.
The highest value in each metric column is bold and the
second-highest is underlined. A double dash denotes an unreported
result. All metrics are higher-is-better; local results report the mean and
population standard deviation across five seeds (42--46). The local TUEV
release contains additional training and validation segments, while retaining
the same held-out test set used by the references. Table~\ref{tab:splits} lists
the resulting split sizes and partition rules for all 12 datasets.

\begin{table*}[t]
\centering
\caption{Split sizes and partition rules for the 12 downstream datasets, as
produced by the local loaders. All partitions are inherited unchanged from the
fixed split code of the \cbramod{} and \reve{} benchmarks. TUEV uses v2.0.1
training and validation counts with the unchanged official 29,421-sample test
set.}
\label{tab:splits}
\scriptsize
\setlength{\tabcolsep}{4pt}
\renewcommand{\arraystretch}{1.06}
\begin{tabularx}{\textwidth}{@{}l l r r r X@{}}
\toprule
Dataset & Unit & Train & Val & Test & Partition (inherited) \\
\midrule
CHB-MIT & segments & 286{,}961 & 23{,}064 & 16{,}963 & subjects CHB01--20 / CHB21--22 / CHB23--24 \\
TUAB & segments & 297{,}103 & 75{,}407 & 36{,}945 & official train/test subjects; training subjects split 80:20 \\
TUEV & segments & 68{,}445 & 15{,}487 & 29{,}421 & same subject protocol; official test set \\
ISRUC & records & 3{,}559 & 468 & 435 & subjects 1--80 / 81--90 / 91--100 \\
FACED & segments & 6{,}720 & 1{,}680 & 1{,}932 & subjects 1--80 / 81--100 / 101--123 \\
SEED-V & trials & 34{,}432 & 42{,}960 & 40{,}352 & per session, trials 1--5 / 6--10 / 11--15 \\
PhysioNet-MI & segments & 6{,}300 & 1{,}734 & 1{,}803 & subjects 1--70 / 71--89 / 90--109 \\
SHU-MI & segments & 7{,}210 & 2{,}431 & 2{,}347 & subjects 1--15 / 16--20 / 21--25 \\
BCIC2020-3 & trials & 4{,}500 & 750 & 750 & official 60/10/10 trials per class \\
Mumtaz2016 & segments & 4{,}891 & 1{,}041 & 1{,}211 & 22/5/6 MDD and 21/4/5 control subjects \\
MentalArithmetic & segments & 1{,}343 & 172 & 192 & subjects 1--28 / 29--32 / 33--36 \\
HMC & epochs & 91{,}248 & 22{,}124 & 23{,}871 & recordings 1--100 / 101--125 / 126--151 (sorted) \\
\bottomrule
\end{tabularx}
\tablenote{CHB-MIT counts exclude five corrupt segments that the loader skips.
ISRUC records contain 20 consecutive 30-second epochs, and metrics are computed
over epochs.}
\end{table*}
\FloatBarrier

\newcommand{\detailhead}[3]{\begin{detailresults}\toprule \textbf{Method} & \textbf{#1} & \textbf{#2} & \textbf{#3} \\
\midrule}
\newcommand{\detailend}{\bottomrule\end{detailresults}}

\subsection{CHB-MIT: Seizure Detection}
CHB-MIT evaluates binary seizure detection in clinical EEG recordings \citep{shoeb2009application}. Each local input contains 16 channels and 2,000 time samples. PR-AUC complements balanced accuracy and ROC-AUC by emphasizing detection of the positive seizure class.
CHB-MIT uses the CBraMod subject split CHB01--20/CHB21--22/CHB23--24. The local row
uses the locked five-seed recipe.
\begin{table}[t]\centering\caption{CHB-MIT under the shared subject split.}\label{tab:detail-chb}
\detailhead{Balanced accuracy}{PR-AUC}{ROC-AUC}
\multicolumn{4}{@{}l}{\textit{Supervised EEG models}} \\
EEGNet& $0.5658\pm0.0106$ & $0.1914\pm0.0182$ & $0.8048\pm0.0136$  \\
EEGConformer& $0.5976\pm0.0141$ & $0.2209\pm0.0215$ & $0.8226\pm0.0170$  \\
SPaRCNet& $0.5876\pm0.0191$ & $0.1247\pm0.0119$ & $0.8143\pm0.0148$  \\
ContraWR& $0.6344\pm0.0002$ & $0.2264\pm0.0174$ & $0.8097\pm0.0114$  \\
CNN-Transformer& $0.6389\pm0.0067$ & $0.2479\pm0.0227$ & $0.8662\pm0.0082$  \\
FFCL& $0.6262\pm0.0104$ & $0.2049\pm0.0346$ & $0.8271\pm0.0051$  \\
ST-Transformer& $0.5915\pm0.0195$ & $0.1422\pm0.0094$ & $0.8237\pm0.0491$  \\
\midrule
\multicolumn{4}{@{}l}{\textit{EEG foundation models}} \\
\biot{}& $0.7068\pm0.0457$ & $0.3277\pm0.0460$ & $0.8761\pm0.0284$  \\
\labram{}-Base& $0.7075\pm0.0358$ & $0.3287\pm0.0402$ & $0.8679\pm0.0199$  \\
\cbramod{}& $0.7398\pm0.0284$ & $0.3689\pm0.0382$ & $0.8892\pm0.0154$  \\
\reve{}-Base& -- & -- & --  \\
\midrule
\multicolumn{4}{@{}l}{\textit{Vision models}} \\
\method{} (EfficientNet-B0)& \secondcell{0.7552\pm0.0599} & $0.4092\pm0.0536$ & $0.8974\pm0.0194$  \\
\method{} (ConvNeXt-Tiny)& \bestcell{0.7682\pm0.0294} & \secondcell{0.4734\pm0.0345} & \secondcell{0.9067\pm0.0156}  \\
\method{} (ViT-Small)& \secondcell{0.7204\pm0.0175} & \bestcell{0.4759\pm0.0118} & \bestcell{0.9116\pm0.0071}  \\
\detailend\end{table}

\subsection{TUAB: Abnormal EEG Detection}
TUAB evaluates binary classification of normal and abnormal clinical EEG \citep{obeid2016temple}. The local pipeline represents each segment by 16 channels and 2,000 time samples. We report balanced accuracy, PR-AUC, and ROC-AUC; the local split contains 297,103 training, 75,407 validation, and 36,945 test segments.
TUAB uses the target-corpus-excluded \cbramod{} control because TUAB is part
of the TUEG pretraining source.
\begin{table}[t]\centering\caption{TUAB under the shared split.}\label{tab:detail-tuab}
\detailhead{Balanced accuracy}{PR-AUC}{ROC-AUC}
\multicolumn{4}{@{}l}{\textit{Supervised EEG models}} \\
EEGNet& $0.7642\pm0.0036$ & $0.8299\pm0.0043$ & $0.8412\pm0.0031$  \\
EEGConformer& $0.7758\pm0.0049$ & $0.8427\pm0.0054$ & $0.8445\pm0.0038$  \\
SPaRCNet& $0.7896\pm0.0018$ & $0.8414\pm0.0018$ & $0.8676\pm0.0012$  \\
ContraWR& $0.7746\pm0.0041$ & $0.8421\pm0.0104$ & $0.8456\pm0.0074$  \\
CNN-Transformer& $0.7777\pm0.0022$ & $0.8433\pm0.0039$ & $0.8461\pm0.0013$  \\
FFCL& $0.7848\pm0.0038$ & $0.8448\pm0.0065$ & $0.8569\pm0.0051$  \\
ST-Transformer& $0.7966\pm0.0023$ & $0.8521\pm0.0026$ & $0.8707\pm0.0019$  \\
\midrule
\multicolumn{4}{@{}l}{\textit{EEG foundation models}} \\
\biot{}& $0.7959\pm0.0057$ & $0.8792\pm0.0023$ & $0.8815\pm0.0043$  \\
\labram{}-Base& $0.8140\pm0.0019$ & $0.8965\pm0.0016$ & $0.9022\pm0.0009$  \\
\cbramod{}& $0.8249\pm0.0025$ & \secondcell{0.9221\pm0.0015} & \secondcell{0.9156\pm0.0017}  \\
\reve{}-Base& \secondcell{0.8315\pm0.0014} & \bestcell{0.9281\pm0.0009} & \bestcell{0.9245\pm0.0013}  \\
\midrule
\multicolumn{4}{@{}l}{\textit{Vision models}} \\
\method{} (EfficientNet-B0)& \bestcell{0.8317\pm0.0033} & $0.9178\pm0.0025$ & $0.9126\pm0.0024$  \\
\method{} (ConvNeXt-Tiny)& $0.8216\pm0.0050$ & $0.9059\pm0.0045$ & $0.9020\pm0.0045$  \\
\method{} (ViT-Small)& $0.8202\pm0.0031$ & $0.9091\pm0.0034$ & $0.9042\pm0.0024$  \\
\detailend\end{table}

\subsection{TUEV: Event Type Classification}
TUEV evaluates six-class EEG event classification \citep{obeid2016temple}. Local inputs contain 16 channels and 1,000 time samples, corresponding to five seconds at 200\,Hz. We report balanced accuracy, Cohen's $\kappa$, and weighted F1 over 68,445 training, 15,487 validation, and 29,421 test segments.
The local files use TUEV v2.0.1 rather than the v2.0.0 release used by the
reference benchmarks; the held-out test partition is the official evaluation
set, unchanged at 29,421 segments, while v2.0.1 changes only the training and
validation segments (68,445/15,487 rather than 67,436/15,634).
\begin{table}[t]\centering\caption{TUEV event classification.}\label{tab:detail-tuev}
\detailhead{Balanced accuracy}{Cohen's $\kappa$}{Weighted F1}
\multicolumn{4}{@{}l}{\textit{Supervised EEG models}} \\
EEGNet& $0.3876\pm0.0143$ & $0.3577\pm0.0155$ & $0.6539\pm0.0120$  \\
EEGConformer& $0.4074\pm0.0164$ & $0.3967\pm0.0195$ & $0.6983\pm0.0152$  \\
SPaRCNet& $0.4161\pm0.0262$ & $0.4233\pm0.0181$ & $0.7024\pm0.0104$  \\
ContraWR& $0.4384\pm0.0349$ & $0.3912\pm0.0237$ & $0.6893\pm0.0136$  \\
CNN-Transformer& $0.4087\pm0.0161$ & $0.3815\pm0.0134$ & $0.6854\pm0.0293$  \\
FFCL& $0.3979\pm0.0104$ & $0.3732\pm0.0188$ & $0.6783\pm0.0120$  \\
ST-Transformer& $0.3984\pm0.0228$ & $0.3765\pm0.0306$ & $0.6823\pm0.0190$  \\
\midrule
\multicolumn{4}{@{}l}{\textit{EEG foundation models}} \\
\biot{}& $0.5281\pm0.0225$ & $0.5273\pm0.0249$ & $0.7492\pm0.0082$  \\
\labram{}-Base& $0.6409\pm0.0065$ & $0.6637\pm0.0093$ & $0.8312\pm0.0052$  \\
\cbramod{}& $0.6659\pm0.0124$ & $0.6744\pm0.0121$ & $0.8331\pm0.0071$  \\
\reve{}-Base& $0.6759\pm0.0229$ & $0.6783\pm0.0199$ & $0.8451\pm0.0129$  \\
\midrule
\multicolumn{4}{@{}l}{\textit{Vision models}} \\
\method{} (EfficientNet-B0)& $0.6423\pm0.0164$ & $0.6901\pm0.0190$ & $0.8380\pm0.0076$  \\
\method{} (ConvNeXt-Tiny)& \secondcell{0.6763\pm0.0122} & \secondcell{0.7095\pm0.0305} & \secondcell{0.8488\pm0.0135}  \\
\method{} (ViT-Small)& \bestcell{0.7261\pm0.0135} & \bestcell{0.7428\pm0.0127} & \bestcell{0.8655\pm0.0055}  \\
\detailend\end{table}

\subsection{ISRUC: Sleep Staging}
ISRUC-Sleep provides overnight polysomnography for five-stage sleep classification \citep{khalighi2016isruc}. Each local epoch contains six channels and 6,000 time samples (30 seconds at 200\,Hz). Balanced accuracy, Cohen's $\kappa$, and weighted F1 are computed over epoch labels.
ISRUC uses the published benchmark split (subjects 1--80/81--90/91--100 for
train/validation/test), identical to the \reve{} protocol. The local loader
retains 20 consecutive 30-second epochs per record, but the vision path
reshapes each $[20,C,T]$ sequence into 20 independent samples before folding.
The backbone therefore sees one 30-second epoch at a time and does not
aggregate features or predictions across the 20-epoch context window. Its
outputs are reshaped back to $[20,K]$ so that the loss and reported metrics are
computed against all 20 epoch labels. Thus the 20-epoch storage format does not
give the visual models an additional temporal-context input.
\begin{table}[t]\centering\caption{ISRUC five-class sleep staging.}\label{tab:detail-isruc}
\detailhead{Balanced accuracy}{Cohen's $\kappa$}{Weighted F1}
\multicolumn{4}{@{}l}{\textit{Supervised EEG models}} \\
EEGNet& $0.7154\pm0.0121$ & $0.7040\pm0.0173$ & $0.7513\pm0.0124$  \\
EEGConformer& $0.7400\pm0.0133$ & $0.7143\pm0.0162$ & $0.7634\pm0.0151$  \\
SPaRCNet& $0.7487\pm0.0075$ & $0.7097\pm0.0132$ & $0.7624\pm0.0092$  \\
ContraWR& $0.7402\pm0.0126$ & $0.7178\pm0.0156$ & $0.7610\pm0.0137$  \\
CNN-Transformer& $0.7363\pm0.0087$ & $0.7129\pm0.0121$ & $0.7719\pm0.0105$  \\
FFCL& $0.7277\pm0.0182$ & $0.7016\pm0.0291$ & $0.7614\pm0.0197$  \\
ST-Transformer& $0.7381\pm0.0205$ & $0.7013\pm0.0352$ & $0.7681\pm0.0175$  \\
\midrule
\multicolumn{4}{@{}l}{\textit{EEG foundation models}} \\
\biot{}& $0.7527\pm0.0121$ & $0.7192\pm0.0231$ & $0.7790\pm0.0146$  \\
\labram{}-Base& $0.7633\pm0.0102$ & $0.7231\pm0.0182$ & $0.7810\pm0.0133$  \\
\cbramod{}& $0.7865\pm0.0110$ & $0.7442\pm0.0152$ & $0.8011\pm0.0099$  \\
\reve{}-Base& $0.7819\pm0.0078$ & $0.7500\pm0.0156$ & $0.8005\pm0.0135$  \\
\midrule
\multicolumn{4}{@{}l}{\textit{Vision models}} \\
\method{} (EfficientNet-B0)& $0.8025\pm0.0027$ & $0.7705\pm0.0032$ & $0.8197\pm0.0029$  \\
\method{} (ConvNeXt-Tiny)& \bestcell{0.8142\pm0.0009} & \bestcell{0.7906\pm0.0018} & \bestcell{0.8372\pm0.0012}  \\
\method{} (ViT-Small)& \secondcell{0.8130\pm0.0027} & \secondcell{0.7875\pm0.0026} & \secondcell{0.8349\pm0.0023}  \\
\detailend\end{table}

\subsection{FACED: Emotion Recognition}
FACED evaluates nine-class emotion recognition using recordings from 123 participants \citep{chen2023large}. The local inputs contain 32 channels and 2,000 time samples, corresponding to 10-second segments at 200\,Hz. The split contains 6,720 training, 1,680 validation, and 1,932 test segments; evaluation uses balanced accuracy, Cohen's $\kappa$, and weighted F1.
FACED uses the CBraMod partition of subjects 1--80/81--100/101--123 for
train/validation/test.
\begin{table}[t]\centering\caption{FACED nine-class emotion recognition.}\label{tab:detail-faced}
\detailhead{Balanced accuracy}{Cohen's $\kappa$}{Weighted F1}
\multicolumn{4}{@{}l}{\textit{Supervised EEG models}} \\
EEGNet& $0.4090\pm0.0122$ & $0.3342\pm0.0251$ & $0.4124\pm0.0141$  \\
EEGConformer& $0.4559\pm0.0125$ & $0.3858\pm0.0186$ & $0.4514\pm0.0107$  \\
SPaRCNet& $0.4673\pm0.0155$ & $0.3978\pm0.0289$ & $0.4729\pm0.0133$  \\
ContraWR& $0.4887\pm0.0078$ & $0.4231\pm0.0151$ & $0.4884\pm0.0074$  \\
CNN-Transformer& $0.4697\pm0.0132$ & $0.4017\pm0.0168$ & $0.4720\pm0.0125$  \\
FFCL& $0.4673\pm0.0158$ & $0.3987\pm0.0383$ & $0.4699\pm0.0145$  \\
ST-Transformer& $0.4810\pm0.0079$ & $0.4137\pm0.0133$ & $0.4795\pm0.0096$  \\
\midrule
\multicolumn{4}{@{}l}{\textit{EEG foundation models}} \\
\biot{}& $0.5118\pm0.0118$ & $0.4476\pm0.0254$ & $0.5136\pm0.0112$  \\
\labram{}-Base& $0.5273\pm0.0107$ & $0.4698\pm0.0188$ & $0.5288\pm0.0102$  \\
\cbramod{}& $0.5509\pm0.0089$ & $0.5041\pm0.0122$ & $0.5618\pm0.0093$  \\
\reve{}-Base& \secondcell{0.5646\pm0.0164} & \secondcell{0.5080\pm0.0191} & \secondcell{0.5659\pm0.0172}  \\
\midrule
\multicolumn{4}{@{}l}{\textit{Vision models}} \\
\method{} (EfficientNet-B0)& $0.5516\pm0.0179$ & $0.4953\pm0.0197$ & $0.5567\pm0.0163$  \\
\method{} (ConvNeXt-Tiny)& \bestcell{0.6363\pm0.0027} & \bestcell{0.5891\pm0.0032} & \bestcell{0.6374\pm0.0030}  \\
\method{} (ViT-Small)& $0.4464\pm0.0102$ & $0.3744\pm0.0109$ & $0.4445\pm0.0098$  \\
\detailend\end{table}

\subsection{SEED-V: Emotion Recognition}
SEED-V evaluates five-class emotion recognition \citep{liu2022comparing}. Each local sample contains 62 channels and 200 time samples. The trial-based protocol yields 34,432 training, 42,960 validation, and 40,352 test samples. We report balanced accuracy, Cohen's $\kappa$, and weighted F1.
SEED-V follows the CBraMod trial-level split: for each session, the first five
trials form training, the middle five validation, and the last five test.
\reve{} does not report this dataset.
\begin{table}[t]\centering\caption{SEED-V five-class emotion recognition.}\label{tab:detail-seedv}
\detailhead{Balanced accuracy}{Cohen's $\kappa$}{Weighted F1}
\multicolumn{4}{@{}l}{\textit{Supervised EEG models}} \\
EEGNet& $0.2961\pm0.0102$ & $0.1006\pm0.0143$ & $0.2749\pm0.0098$  \\
EEGConformer& $0.3537\pm0.0112$ & $0.1772\pm0.0174$ & $0.3487\pm0.0136$  \\
SPaRCNet& $0.2949\pm0.0078$ & $0.1121\pm0.0139$ & $0.2979\pm0.0083$  \\
ContraWR& $0.3546\pm0.0105$ & $0.1905\pm0.0188$ & $0.3544\pm0.0121$  \\
CNN-Transformer& $0.3678\pm0.0078$ & $0.2072\pm0.0183$ & $0.3642\pm0.0088$  \\
FFCL& $0.3641\pm0.0092$ & $0.2078\pm0.0201$ & $0.3645\pm0.0132$  \\
ST-Transformer& $0.3052\pm0.0072$ & $0.1083\pm0.0121$ & $0.2833\pm0.0105$  \\
\midrule
\multicolumn{4}{@{}l}{\textit{EEG foundation models}} \\
\biot{}& $0.3837\pm0.0187$ & $0.2261\pm0.0262$ & $0.3856\pm0.0203$  \\
\labram{}-Base& \secondcell{0.3976\pm0.0138} & $0.2386\pm0.0209$ & $0.3974\pm0.0111$  \\
\cbramod{}& \bestcell{0.4091\pm0.0097} & \bestcell{0.2569\pm0.0143} & \bestcell{0.4101\pm0.0108}  \\
\reve{}-Base& -- & -- & --  \\
\midrule
\multicolumn{4}{@{}l}{\textit{Vision models}} \\
\method{} (EfficientNet-B0)& $0.3660\pm0.0062$ & $0.2056\pm0.0078$ & $0.3678\pm0.0057$  \\
\method{} (ConvNeXt-Tiny)& $0.3884\pm0.0071$ & $0.2396\pm0.0095$ & $0.3976\pm0.0081$  \\
\method{} (ViT-Small)& $0.3930\pm0.0035$ & \secondcell{0.2442\pm0.0045} & \secondcell{0.4016\pm0.0038}  \\
\detailend\end{table}

\subsection{PhysioNet-MI: Motor Imagery}
PhysioNet-MI evaluates four-class motor decoding using the PhysioNet motor movement/imagery recordings \citep{schalk2004bci2000}. Local segments contain 64 channels and 800 time samples. The split yields 6,300 training, 1,734 validation, and 1,803 test samples; evaluation uses balanced accuracy, Cohen's $\kappa$, and weighted F1.
PhysioNet-MI uses the shared CBraMod/REVE subject split 1--70/71--89/90--109
and $P=1$ in the local model.
\begin{table}[t]\centering\caption{PhysioNet-MI four-class motor imagery.}\label{tab:detail-physio}
\detailhead{Balanced accuracy}{Cohen's $\kappa$}{Weighted F1}
\multicolumn{4}{@{}l}{\textit{Supervised EEG models}} \\
EEGNet& $0.5814\pm0.0125$ & $0.4468\pm0.0199$ & $0.5796\pm0.0115$  \\
EEGConformer& $0.6049\pm0.0104$ & $0.4736\pm0.0171$ & $0.6062\pm0.0095$  \\
SPaRCNet& $0.5932\pm0.0152$ & $0.4564\pm0.0234$ & $0.5937\pm0.0147$  \\
ContraWR& $0.5892\pm0.0133$ & $0.4527\pm0.0248$ & $0.5918\pm0.0116$  \\
CNN-Transformer& $0.6053\pm0.0118$ & $0.4725\pm0.0223$ & $0.6041\pm0.0105$  \\
FFCL& $0.5726\pm0.0092$ & $0.4323\pm0.0182$ & $0.5701\pm0.0079$  \\
ST-Transformer& $0.6035\pm0.0081$ & $0.4712\pm0.0199$ & $0.6053\pm0.0075$  \\
\midrule
\multicolumn{4}{@{}l}{\textit{EEG foundation models}} \\
\biot{}& $0.6153\pm0.0154$ & $0.4875\pm0.0272$ & $0.6158\pm0.0197$  \\
\labram{}-Base& $0.6173\pm0.0122$ & $0.4912\pm0.0192$ & $0.6177\pm0.0141$  \\
\cbramod{}& $0.6417\pm0.0091$ & $0.5222\pm0.0169$ & $0.6427\pm0.0100$  \\
\reve{}-Base& $0.6480\pm0.0140$ & $0.5306\pm0.0187$ & $0.6484\pm0.0170$  \\
\midrule
\multicolumn{4}{@{}l}{\textit{Vision models}} \\
\method{} (EfficientNet-B0)& \secondcell{0.6491\pm0.0125} & \secondcell{0.5322\pm0.0167} & \secondcell{0.6485\pm0.0133}  \\
\method{} (ConvNeXt-Tiny)& \bestcell{0.6609\pm0.0079} & \bestcell{0.5479\pm0.0105} & \bestcell{0.6618\pm0.0079}  \\
\method{} (ViT-Small)& $0.6377\pm0.0042$ & $0.5169\pm0.0056$ & $0.6386\pm0.0041$  \\
\detailend\end{table}

\subsection{SHU-MI: Motor Imagery}
SHU-MI evaluates binary motor-imagery classification \citep{ma2022large}. Local segments contain 32 channels and 800 time samples. Subject-disjoint evaluation uses 7,210 training, 2,431 validation, and 2,347 test segments, with balanced accuracy, PR-AUC, and ROC-AUC as metrics.
SHU-MI uses the CBraMod subject split 1--15/16--20/21--25; \reve{} does not
report this dataset.
\begin{table}[t]\centering\caption{SHU-MI binary motor imagery.}\label{tab:detail-shu}
\detailhead{Balanced accuracy}{PR-AUC}{ROC-AUC}
\multicolumn{4}{@{}l}{\textit{Supervised EEG models}} \\
EEGNet& $0.5889\pm0.0177$ & $0.6311\pm0.0142$ & $0.6283\pm0.0152$  \\
EEGConformer& $0.5900\pm0.0107$ & $0.6370\pm0.0093$ & $0.6351\pm0.0101$  \\
SPaRCNet& $0.5978\pm0.0097$ & $0.6510\pm0.0062$ & $0.6431\pm0.0082$  \\
ContraWR& $0.5873\pm0.0128$ & $0.6315\pm0.0105$ & $0.6273\pm0.0113$  \\
CNN-Transformer& $0.5975\pm0.0169$ & $0.6412\pm0.0076$ & $0.6323\pm0.0082$  \\
FFCL& $0.5692\pm0.0252$ & $0.5943\pm0.0172$ & $0.6014\pm0.0168$  \\
ST-Transformer& $0.5992\pm0.0206$ & $0.6394\pm0.0122$ & $0.6431\pm0.0111$  \\
\midrule
\multicolumn{4}{@{}l}{\textit{EEG foundation models}} \\
\biot{}& $0.6179\pm0.0183$ & $0.6770\pm0.0119$ & $0.6609\pm0.0127$  \\
\labram{}-Base& $0.6166\pm0.0192$ & $0.6761\pm0.0083$ & $0.6604\pm0.0091$  \\
\cbramod{}& \bestcell{0.6370\pm0.0151} & \bestcell{0.7139\pm0.0088} & \secondcell{0.6988\pm0.0068}  \\
\reve{}-Base& -- & -- & --  \\
\midrule
\multicolumn{4}{@{}l}{\textit{Vision models}} \\
\method{} (EfficientNet-B0)& $0.6273\pm0.0085$ & \secondcell{0.6982\pm0.0178} & $0.6962\pm0.0139$  \\
\method{} (ConvNeXt-Tiny)& \secondcell{0.6358\pm0.0075} & $0.6910\pm0.0129$ & \bestcell{0.7006\pm0.0112}  \\
\method{} (ViT-Small)& $0.6302\pm0.0331$ & $0.6891\pm0.0457$ & $0.6908\pm0.0452$  \\
\detailend\end{table}

\subsection{BCIC2020-3: Imagined Speech}
BCIC2020-3 evaluates five-class imagined-speech decoding \citep{jeong20222020}. Each local input contains 64 channels and 600 time samples. The benchmark contains 4,500 training, 750 validation, and 750 test segments. We report balanced accuracy, Cohen's $\kappa$, and weighted F1.
BCIC2020-3 uses the official 60/10/10 trials-per-class split, the same
partition used by the CBraMod and REVE benchmarks.
\begin{table}[t]\centering\caption{BCIC2020-3 five-class imagined speech.}\label{tab:detail-bcic}
\detailhead{Balanced accuracy}{Cohen's $\kappa$}{Weighted F1}
\multicolumn{4}{@{}l}{\textit{Supervised EEG models}} \\
EEGNet& $0.4413\pm0.0096$ & $0.3016\pm0.0123$ & $0.4413\pm0.0102$  \\
EEGConformer& $0.4506\pm0.0133$ & $0.3133\pm0.0183$ & $0.4488\pm0.0154$  \\
SPaRCNet& $0.4426\pm0.0156$ & $0.3033\pm0.0233$ & $0.4420\pm0.0108$  \\
ContraWR& $0.4257\pm0.0162$ & $0.3078\pm0.0218$ & $0.4407\pm0.0182$  \\
CNN-Transformer& $0.4533\pm0.0092$ & $0.3166\pm0.0118$ & $0.4506\pm0.0127$  \\
FFCL& $0.4678\pm0.0197$ & $0.3301\pm0.0359$ & $0.4689\pm0.0205$  \\
ST-Transformer& $0.4126\pm0.0122$ & $0.2941\pm0.0159$ & $0.4247\pm0.0138$  \\
\midrule
\multicolumn{4}{@{}l}{\textit{EEG foundation models}} \\
\biot{}& $0.4920\pm0.0086$ & $0.3650\pm0.0176$ & $0.4917\pm0.0079$  \\
\labram{}-Base& $0.5060\pm0.0155$ & $0.3800\pm0.0242$ & $0.5054\pm0.0205$  \\
\cbramod{}& $0.5373\pm0.0108$ & $0.4216\pm0.0163$ & $0.5383\pm0.0096$  \\
\reve{}-Base& $0.5635\pm0.0123$ & $0.4543\pm0.0154$ & $0.5633\pm0.0124$  \\
\midrule
\multicolumn{4}{@{}l}{\textit{Vision models}} \\
\method{} (EfficientNet-B0)& \secondcell{0.6539\pm0.0128} & \secondcell{0.5673\pm0.0160} & \secondcell{0.6539\pm0.0130}  \\
\method{} (ConvNeXt-Tiny)& $0.6448\pm0.0091$ & $0.5560\pm0.0114$ & $0.6447\pm0.0092$  \\
\method{} (ViT-Small)& \bestcell{0.6619\pm0.0090} & \bestcell{0.5773\pm0.0113} & \bestcell{0.6618\pm0.0089}  \\
\detailend\end{table}

\subsection{Mumtaz2016: Depression Detection}
Mumtaz2016 evaluates binary discrimination between major depressive disorder (MDD) and healthy controls \citep{mumtaz2016mdd}. Each local input contains 19 channels and 1,000 time samples. The fixed benchmark partition is subject-disjoint and yields 4,891 training, 1,041 validation, and 1,211 test segments; evaluation uses balanced accuracy, PR-AUC, and ROC-AUC.
The partition assigns 22/5/6 MDD and 21/4/5 control subjects to
train/validation/test (63 subjects; both eyes-open and eyes-closed recordings),
matching the shared CBraMod/REVE preprocessing code.
\begin{table}[t]\centering\caption{Mumtaz2016 depression detection.}\label{tab:detail-mumtaz}
\detailhead{Balanced accuracy}{PR-AUC}{ROC-AUC}
\multicolumn{4}{@{}l}{\textit{Supervised EEG models}} \\
EEGNet& $0.9232\pm0.0104$ & $0.9626\pm0.0095$ & $0.9639\pm0.0093$  \\
EEGConformer& $0.9308\pm0.0117$ & $0.9684\pm0.0105$ & $0.9702\pm0.0101$  \\
SPaRCNet& $0.9316\pm0.0095$ & $0.9754\pm0.0065$ & $0.9781\pm0.0083$  \\
ContraWR& $0.9195\pm0.0115$ & $0.9589\pm0.0102$ & $0.9621\pm0.0092$  \\
CNN-Transformer& $0.9305\pm0.0068$ & $0.9757\pm0.0074$ & $0.9742\pm0.0059$  \\
FFCL& $0.9314\pm0.0038$ & $0.9717\pm0.0021$ & $0.9753\pm0.0033$  \\
ST-Transformer& $0.9135\pm0.0103$ & $0.9578\pm0.0086$ & $0.9594\pm0.0059$  \\
\midrule
\multicolumn{4}{@{}l}{\textit{EEG foundation models}} \\
\biot{}& $0.9358\pm0.0052$ & $0.9736\pm0.0034$ & $0.9758\pm0.0042$  \\
\labram{}-Base& $0.9409\pm0.0079$ & $0.9798\pm0.0093$ & $0.9782\pm0.0057$  \\
\cbramod{}& \secondcell{0.9560\pm0.0056} & \secondcell{0.9923\pm0.0032} & \secondcell{0.9921\pm0.0025}  \\
\reve{}-Base& \bestcell{0.9644\pm0.0097} & \bestcell{0.9961\pm0.0013} & \bestcell{0.9957\pm0.0015}  \\
\midrule
\multicolumn{4}{@{}l}{\textit{Vision models}} \\
\method{} (EfficientNet-B0)& $0.9153\pm0.0061$ & $0.9864\pm0.0045$ & $0.9846\pm0.0053$  \\
\method{} (ConvNeXt-Tiny)& $0.9176\pm0.0138$ & $0.9842\pm0.0020$ & $0.9825\pm0.0034$  \\
\method{} (ViT-Small)& $0.9112\pm0.0099$ & $0.9797\pm0.0037$ & $0.9767\pm0.0047$  \\
\detailend\end{table}

\subsection{MentalArithmetic: Mental Workload}
MentalArithmetic evaluates binary workload classification from EEG recorded before and during mental arithmetic \citep{zyma2019mental}. Local inputs contain 20 channels and 1,000 time samples. The subject-disjoint split yields 1,343 training, 172 validation, and 192 test segments. We report balanced accuracy, PR-AUC, and ROC-AUC.
MentalArithmetic uses the shared CBraMod/REVE subject split
1--28/29--32/33--36 for train/validation/test.
\begin{table}[t]\centering\caption{MentalArithmetic binary workload classification.}\label{tab:detail-mental}
\detailhead{Balanced accuracy}{PR-AUC}{ROC-AUC}
\multicolumn{4}{@{}l}{\textit{Supervised EEG models}} \\
EEGNet& $0.6770\pm0.0116$ & $0.5763\pm0.0102$ & $0.7321\pm0.0108$  \\
EEGConformer& $0.6805\pm0.0123$ & $0.5829\pm0.0134$ & $0.7424\pm0.0128$  \\
SPaRCNet& $0.6879\pm0.0107$ & $0.5825\pm0.0193$ & $0.7418\pm0.0132$  \\
ContraWR& $0.6631\pm0.0097$ & $0.5787\pm0.0164$ & $0.7332\pm0.0082$  \\
CNN-Transformer& $0.6779\pm0.0268$ & $0.5777\pm0.0285$ & $0.7258\pm0.0336$  \\
FFCL& $0.6798\pm0.0142$ & $0.5786\pm0.0266$ & $0.7330\pm0.0198$  \\
ST-Transformer& $0.6631\pm0.0173$ & $0.5672\pm0.0259$ & $0.7132\pm0.0174$  \\
\midrule
\multicolumn{4}{@{}l}{\textit{EEG foundation models}} \\
\biot{}& $0.6875\pm0.0186$ & $0.6004\pm0.0195$ & $0.7536\pm0.0144$  \\
\labram{}-Base& $0.6909\pm0.0125$ & $0.5999\pm0.0155$ & $0.7721\pm0.0093$  \\
\cbramod{}& $0.7256\pm0.0132$ & $0.6267\pm0.0099$ & $0.7905\pm0.0073$  \\
\reve{}-Base& \secondcell{0.7660\pm0.0355} & $0.7470\pm0.0807$ & $0.8450\pm0.0514$  \\
\midrule
\multicolumn{4}{@{}l}{\textit{Vision models}} \\
\method{} (EfficientNet-B0)& $0.6910\pm0.0418$ & \bestcell{0.7903\pm0.0536} & \bestcell{0.8779\pm0.0251}  \\
\method{} (ConvNeXt-Tiny)& \bestcell{0.7667\pm0.0234} & \secondcell{0.7598\pm0.0436} & \secondcell{0.8752\pm0.0313}  \\
\method{} (ViT-Small)& $0.7507\pm0.0464$ & $0.7249\pm0.0625$ & $0.8498\pm0.0433$  \\
\detailend\end{table}


\subsection{HMC: Sleep Staging}
HMC evaluates five-stage sleep classification from polysomnography \citep{alvarez2022hmc}. Each local epoch contains four channels and 6,000 time samples. The split yields 91,248 training, 22,124 validation, and 23,871 test epochs (137,243 in total); evaluation uses balanced accuracy, Cohen's $\kappa$, and weighted F1.
HMC follows the \reve{} benchmark preprocessing and partition, itself adapted
from NeuroLM: the 151 recordings are sorted by filename and assigned 100/25/26
to train, validation, and test, yielding four-channel, 30-second PSG epochs.
The local configuration uses $P=16$ temporal folding,
validation-selected checkpoints, and one final test evaluation per seed.
Published rows are transcribed from \reve{}'s sleep-staging table, which uses
the same held-out recordings; \cbramod{} does not report this corpus.
\begin{table}[t]\centering\caption{HMC five-class sleep staging.}\label{tab:detail-hmc}
\detailhead{Balanced accuracy}{Cohen's $\kappa$}{Weighted F1}
\multicolumn{4}{@{}l}{\textit{Supervised EEG models}} \\
SPaRCNet& $0.4756\pm0.1109$ & $0.3147\pm0.1315$ & $0.4108\pm0.1310$  \\
ContraWR& $0.4242\pm0.0541$ & $0.2340\pm0.0554$ & $0.2987\pm0.0288$  \\
CNN-Transformer& $0.6573\pm0.0141$ & $0.5961\pm0.0105$ & $0.6896\pm0.0065$  \\
FFCL& $0.4427\pm0.0702$ & $0.2542\pm0.0654$ & $0.2902\pm0.0485$  \\
ST-Transformer& $0.2559\pm0.0141$ & $0.0503\pm0.0183$ & $0.1428\pm0.0122$  \\
\midrule
\multicolumn{4}{@{}l}{\textit{EEG foundation models}} \\
\biot{}& $0.6862\pm0.0041$ & $0.6295\pm0.0113$ & $0.7091\pm0.0147$  \\
\labram{}-Base& $0.7286\pm0.0101$ & $0.6812\pm0.0073$ & $0.7554\pm0.0024$  \\
\reve{}-Base& $0.7401\pm0.0075$ & $0.6982\pm0.0078$ & $0.7638\pm0.0074$  \\
\midrule
\multicolumn{4}{@{}l}{\textit{Vision models}} \\
\method{} (EfficientNet-B0) & \secondcell{0.7557\pm0.0019} & \secondcell{0.7070\pm0.0023} & \secondcell{0.7700\pm0.0017} \\
\method{} (ConvNeXt-Tiny) & \bestcell{0.7581\pm0.0055} & \bestcell{0.7111\pm0.0063} & \bestcell{0.7744\pm0.0043} \\
\method{} (ViT-Small) & $0.7476\pm0.0032$ & $0.6987\pm0.0033$ & $0.7634\pm0.0017$ \\
\detailend\end{table}


\section{Matched Foundation-Model Reruns}
\label{app:foundation-rerun}

We reran the released CBraMod and REVE-Base checkpoints on TUAB, FACED, and
HMC using the inherited partitions, seeds 42--46, validation-selected
checkpoints, and one final test evaluation per seed. Table~\ref{tab:foundation-reruns}
lists these reruns next to our finalized vision models, which follow the
identical protocol, so that the two families can be compared under one
pipeline. The downstream pipeline
and optimizer match our local experiments, while each model retains its
released input normalization. These are matched-protocol reruns rather than
exact reproductions of every original training stage.

CBraMod's released multi-learning-rate recipe is replaced by one uniform
learning rate. REVE's published linear-probe warm start, mixup, StableAdamW,
LoRA, and model souping are not applied. For TUAB, REVE receives the same
16-channel bipolar arrays used by CBraMod rather than the released 21-channel
memmap. These deviations can lower rerun values relative to published results,
so the experiment quantifies protocol sensitivity rather than replacing the
original references.

\begin{table*}[t]
\centering\scriptsize
\caption{Protocol-matched comparison on the three rerun datasets. Released
CBraMod and REVE-Base checkpoints are rerun under our downstream pipeline
(inherited partitions, seeds 42--46, validation-selected checkpoints, one final
test per seed) and are listed together with our finalized vision models, which
follow the same protocol. Values are five-seed means $\pm$ population standard
deviations; the primary metric is the validation selection metric (PR-AUC for
TUAB, $\kappa$ for FACED and HMC). Bold marks the best matched-protocol value
in each column; published values are transcribed from the original benchmark
tables and are not ranked.}
\label{tab:foundation-reruns}
\setlength{\tabcolsep}{6pt}
\renewcommand{\arraystretch}{1.08}
\begin{tabular}{@{}llccc@{}}
\toprule
Dataset & Method & BA & Primary metric & Published \\
\midrule
\multirow{5}{*}{TUAB}
  & CBraMod (rerun) & $0.7821\pm0.0099$ & $0.8632\pm0.0056$ (PR-AUC) & $0.9221$ \\
  & REVE-Base (rerun) & $0.8052\pm0.0072$ & $0.8901\pm0.0114$ (PR-AUC) & $0.9281$ \\
\cmidrule(lr){2-5}
  & EfficientNet-B0 & \bestcell{0.8317\pm0.0033} & \bestcell{0.9178\pm0.0025} (PR-AUC) & -- \\
  & ConvNeXt-Tiny & $0.8216\pm0.0050$ & $0.9059\pm0.0045$ (PR-AUC) & -- \\
  & ViT-Small & $0.8202\pm0.0031$ & $0.9091\pm0.0034$ (PR-AUC) & -- \\
\midrule
\multirow{5}{*}{FACED}
  & CBraMod (rerun) & $0.5551\pm0.0107$ & $0.4961\pm0.0124$ ($\kappa$) & $0.5041$ \\
  & REVE-Base (rerun) & $0.4850\pm0.0062$ & $0.4190\pm0.0063$ ($\kappa$) & $0.5080$ \\
\cmidrule(lr){2-5}
  & EfficientNet-B0 & $0.5516\pm0.0179$ & $0.4953\pm0.0197$ ($\kappa$) & -- \\
  & ConvNeXt-Tiny & \bestcell{0.6363\pm0.0027} & \bestcell{0.5891\pm0.0032} ($\kappa$) & -- \\
  & ViT-Small & $0.4464\pm0.0102$ & $0.3744\pm0.0109$ ($\kappa$) & -- \\
\midrule
\multirow{5}{*}{HMC}
  & CBraMod (rerun) & $0.7279\pm0.0038$ & $0.6758\pm0.0026$ ($\kappa$) & -- \\
  & REVE-Base (rerun) & $0.7337\pm0.0031$ & $0.6831\pm0.0046$ ($\kappa$) & $0.6982$ \\
\cmidrule(lr){2-5}
  & EfficientNet-B0 & $0.7557\pm0.0019$ & $0.7070\pm0.0023$ ($\kappa$) & -- \\
  & ConvNeXt-Tiny & \bestcell{0.7581\pm0.0055} & \bestcell{0.7111\pm0.0063} ($\kappa$) & -- \\
  & ViT-Small & $0.7476\pm0.0032$ & $0.6987\pm0.0033$ ($\kappa$) & -- \\
\bottomrule
\end{tabular}
\tablenote{Published values are transcribed from the CBraMod and REVE
benchmark tables; HMC is not reported by CBraMod. The reruns use a common
downstream protocol and are not strict replications of all original training
stages.}
\end{table*}

Under this matched protocol the two foundation models fall below their
published values, which quantifies how much those numbers depend on the
original training recipes. Our vision backbones are competitive under the same
protocol: EfficientNet-B0 leads all five methods on TUAB in both balanced
accuracy and PR-AUC, and ConvNeXt-Tiny leads on FACED and HMC. The only near
ties are on FACED, where EfficientNet-B0 matches the rerun CBraMod in both
balanced accuracy ($0.5516$ versus $0.5551$) and $\kappa$ ($0.4953$ versus
$0.4961$), whereas the rerun REVE-Base is clearly lower there than its
published value. We therefore read this table as a protocol-controlled
comparison and keep the published values as the external reference in the main
table.

\clearpage
\section{Ablation Studies}
\label{app:ablations}
\subsection{Matched visual-pretraining initialization}\label{app:pretraining-full-generated}\label{app:complete-ablations}
This section reports the four ablation modules in the order used in the main
text: matched visual initialization, frozen linear probing, temporal folding
versus contiguous reshape, and fold-factor sensitivity. Initialization controls
change backbone initialization; layout controls change the temporal-to-row
mapping. All values are final test means $\pm$
population standard deviations over five seeds; each $\Delta$ is computed as
the first condition minus the second and is reported in percentage points.


\begin{table}[tbp]\centering\scriptsize\caption{Matched initialization (efficientnet-b0).}\label{tab:app-init-effnet}\begin{tabular}{@{}l l c c c@{}}\toprule Dataset & Metric & Pretrained & Random & $\Delta$ \\ \midrule
SEED-V & ba & 0.3660 $\pm$ 0.0062 & 0.3785 $\pm$ 0.0129 & -1.25 \\
SEED-V & kappa & 0.2056 $\pm$ 0.0078 & 0.2266 $\pm$ 0.0170 & -2.10 \\
SEED-V & f1 & 0.3678 $\pm$ 0.0057 & 0.3868 $\pm$ 0.0150 & -1.90 \\
PhysioNet-MI & ba & 0.6491 $\pm$ 0.0125 & 0.5931 $\pm$ 0.0138 & +5.60 \\
PhysioNet-MI & kappa & 0.5322 $\pm$ 0.0167 & 0.4573 $\pm$ 0.0184 & +7.48 \\
PhysioNet-MI & f1 & 0.6485 $\pm$ 0.0133 & 0.5972 $\pm$ 0.0126 & +5.13 \\
Mumtaz2016 & ba & 0.9153 $\pm$ 0.0061 & 0.9175 $\pm$ 0.0115 & -0.23 \\
Mumtaz2016 & pr-auc & 0.9864 $\pm$ 0.0045 & 0.9773 $\pm$ 0.0060 & +0.91 \\
Mumtaz2016 & roc-auc & 0.9846 $\pm$ 0.0053 & 0.9734 $\pm$ 0.0072 & +1.12 \\
ISRUC & ba & 0.8025 $\pm$ 0.0027 & 0.7989 $\pm$ 0.0030 & +0.37 \\
ISRUC & kappa & 0.7704 $\pm$ 0.0032 & 0.7733 $\pm$ 0.0038 & -0.28 \\
ISRUC & f1 & 0.8197 $\pm$ 0.0029 & 0.8243 $\pm$ 0.0029 & -0.46 \\
MentalArithmetic & ba & 0.6910 $\pm$ 0.0418 & 0.5000 $\pm$ 0.0000 & +19.10 \\
MentalArithmetic & pr-auc & 0.7903 $\pm$ 0.0536 & 0.2054 $\pm$ 0.0702 & +58.49 \\
MentalArithmetic & roc-auc & 0.8779 $\pm$ 0.0251 & 0.4603 $\pm$ 0.0323 & +41.76 \\
FACED & ba & 0.5516 $\pm$ 0.0179 & 0.2455 $\pm$ 0.0531 & +30.62 \\
FACED & kappa & 0.4953 $\pm$ 0.0197 & 0.1513 $\pm$ 0.0599 & +34.40 \\
FACED & f1 & 0.5567 $\pm$ 0.0163 & 0.2306 $\pm$ 0.0821 & +32.61 \\
SHU-MI & ba & 0.6273 $\pm$ 0.0085 & 0.5005 $\pm$ 0.0004 & +12.68 \\
SHU-MI & pr-auc & 0.6982 $\pm$ 0.0178 & 0.5157 $\pm$ 0.0175 & +18.25 \\
SHU-MI & roc-auc & 0.6962 $\pm$ 0.0139 & 0.5174 $\pm$ 0.0201 & +17.88 \\
HMC & ba & 0.7557 $\pm$ 0.0019 & 0.7355 $\pm$ 0.0044 & +2.02 \\
HMC & kappa & 0.7070 $\pm$ 0.0023 & 0.6877 $\pm$ 0.0050 & +1.93 \\
HMC & f1 & 0.7700 $\pm$ 0.0017 & 0.7565 $\pm$ 0.0043 & +1.35 \\
TUAB & ba & 0.8317 $\pm$ 0.0033 & 0.8287 $\pm$ 0.0024 & +0.30 \\
TUAB & pr-auc & 0.9178 $\pm$ 0.0025 & 0.9165 $\pm$ 0.0017 & +0.13 \\
TUAB & roc-auc & 0.9126 $\pm$ 0.0024 & 0.9118 $\pm$ 0.0025 & +0.08 \\
TUEV & ba & 0.6423 $\pm$ 0.0164 & 0.5328 $\pm$ 0.0357 & +10.95 \\
TUEV & kappa & 0.6901 $\pm$ 0.0190 & 0.5995 $\pm$ 0.0346 & +9.06 \\
TUEV & f1 & 0.8380 $\pm$ 0.0076 & 0.7914 $\pm$ 0.0196 & +4.66 \\
CHB-MIT & ba & 0.7552 $\pm$ 0.0599 & 0.7300 $\pm$ 0.0453 & +2.52 \\
CHB-MIT & pr-auc & 0.4092 $\pm$ 0.0536 & 0.4037 $\pm$ 0.0224 & +0.55 \\
CHB-MIT & roc-auc & 0.8974 $\pm$ 0.0194 & 0.8882 $\pm$ 0.0220 & +0.92 \\
BCIC2020-3 & ba & 0.6539 $\pm$ 0.0128 & 0.4416 $\pm$ 0.0470 & +21.23 \\
BCIC2020-3 & kappa & 0.5673 $\pm$ 0.0160 & 0.3020 $\pm$ 0.0588 & +26.53 \\
BCIC2020-3 & f1 & 0.6539 $\pm$ 0.0129 & 0.4425 $\pm$ 0.0472 & +21.13 \\
\bottomrule\end{tabular}\end{table}

\begin{table}[tbp]\centering\scriptsize\caption{Matched initialization (convnext-tiny-dinov3).}\label{tab:app-init-convnext}\begin{tabular}{@{}l l c c c@{}}\toprule Dataset & Metric & Pretrained & Random & $\Delta$ \\ \midrule
SEED-V & ba & 0.3884 $\pm$ 0.0071 & 0.2138 $\pm$ 0.0011 & +17.46 \\
SEED-V & kappa & 0.2396 $\pm$ 0.0095 & 0.0191 $\pm$ 0.0016 & +22.05 \\
SEED-V & f1 & 0.3976 $\pm$ 0.0080 & 0.2144 $\pm$ 0.0033 & +18.32 \\
HMC & ba & 0.7581 $\pm$ 0.0055 & 0.7035 $\pm$ 0.0045 & +5.46 \\
HMC & kappa & 0.7111 $\pm$ 0.0063 & 0.6461 $\pm$ 0.0048 & +6.51 \\
HMC & f1 & 0.7744 $\pm$ 0.0043 & 0.7232 $\pm$ 0.0042 & +5.12 \\
SHU-MI & ba & 0.6358 $\pm$ 0.0075 & 0.5775 $\pm$ 0.0093 & +5.83 \\
SHU-MI & pr-auc & 0.6910 $\pm$ 0.0129 & 0.5983 $\pm$ 0.0076 & +9.26 \\
SHU-MI & roc-auc & 0.7006 $\pm$ 0.0112 & 0.6144 $\pm$ 0.0061 & +8.62 \\
PhysioNet-MI & ba & 0.6609 $\pm$ 0.0079 & 0.4523 $\pm$ 0.0089 & +20.86 \\
PhysioNet-MI & kappa & 0.5478 $\pm$ 0.0105 & 0.2697 $\pm$ 0.0119 & +27.82 \\
PhysioNet-MI & f1 & 0.6618 $\pm$ 0.0079 & 0.4529 $\pm$ 0.0088 & +20.89 \\
ISRUC & ba & 0.8142 $\pm$ 0.0009 & 0.6807 $\pm$ 0.0207 & +13.34 \\
ISRUC & kappa & 0.7905 $\pm$ 0.0018 & 0.6378 $\pm$ 0.0216 & +15.27 \\
ISRUC & f1 & 0.8371 $\pm$ 0.0012 & 0.7111 $\pm$ 0.0203 & +12.61 \\
CHB-MIT & ba & 0.7681 $\pm$ 0.0294 & 0.7017 $\pm$ 0.0083 & +6.64 \\
CHB-MIT & pr-auc & 0.4734 $\pm$ 0.0345 & 0.4279 $\pm$ 0.0094 & +4.55 \\
CHB-MIT & roc-auc & 0.9067 $\pm$ 0.0156 & 0.8150 $\pm$ 0.0045 & +9.17 \\
TUAB & ba & 0.8216 $\pm$ 0.0050 & 0.8228 $\pm$ 0.0023 & -0.13 \\
TUAB & pr-auc & 0.9059 $\pm$ 0.0045 & 0.9032 $\pm$ 0.0008 & +0.26 \\
TUAB & roc-auc & 0.9020 $\pm$ 0.0045 & 0.8997 $\pm$ 0.0013 & +0.23 \\
MentalArithmetic & ba & 0.7667 $\pm$ 0.0234 & 0.4875 $\pm$ 0.0316 & +27.92 \\
MentalArithmetic & pr-auc & 0.7598 $\pm$ 0.0436 & 0.2511 $\pm$ 0.0262 & +50.87 \\
MentalArithmetic & roc-auc & 0.8752 $\pm$ 0.0313 & 0.4971 $\pm$ 0.0448 & +37.81 \\
TUEV & ba & 0.6763 $\pm$ 0.0122 & 0.4302 $\pm$ 0.0300 & +24.60 \\
TUEV & kappa & 0.7095 $\pm$ 0.0305 & 0.5070 $\pm$ 0.0094 & +20.25 \\
TUEV & f1 & 0.8488 $\pm$ 0.0135 & 0.7409 $\pm$ 0.0044 & +10.79 \\
Mumtaz2016 & ba & 0.9176 $\pm$ 0.0138 & 0.8987 $\pm$ 0.0061 & +1.89 \\
Mumtaz2016 & pr-auc & 0.9842 $\pm$ 0.0020 & 0.9710 $\pm$ 0.0023 & +1.32 \\
Mumtaz2016 & roc-auc & 0.9825 $\pm$ 0.0034 & 0.9659 $\pm$ 0.0032 & +1.65 \\
BCIC2020-3 & ba & 0.6448 $\pm$ 0.0091 & 0.2947 $\pm$ 0.0187 & +35.01 \\
BCIC2020-3 & kappa & 0.5560 $\pm$ 0.0114 & 0.1183 $\pm$ 0.0234 & +43.77 \\
BCIC2020-3 & f1 & 0.6447 $\pm$ 0.0092 & 0.2939 $\pm$ 0.0191 & +35.08 \\
FACED & ba & 0.6363 $\pm$ 0.0027 & 0.3259 $\pm$ 0.0115 & +31.05 \\
FACED & kappa & 0.5890 $\pm$ 0.0032 & 0.2397 $\pm$ 0.0127 & +34.94 \\
FACED & f1 & 0.6374 $\pm$ 0.0030 & 0.3267 $\pm$ 0.0109 & +31.06 \\
\bottomrule\end{tabular}\end{table}

\begin{table}[tbp]\centering\scriptsize\caption{Matched initialization (vit-small-dinov3).}\label{tab:app-init-vit}\begin{tabular}{@{}l l c c c@{}}\toprule Dataset & Metric & Pretrained & Random & $\Delta$ \\ \midrule
Mumtaz2016 & ba & 0.9112 $\pm$ 0.0099 & 0.9021 $\pm$ 0.0138 & +0.91 \\
Mumtaz2016 & pr-auc & 0.9797 $\pm$ 0.0037 & 0.9778 $\pm$ 0.0020 & +0.19 \\
Mumtaz2016 & roc-auc & 0.9767 $\pm$ 0.0047 & 0.9754 $\pm$ 0.0028 & +0.13 \\
PhysioNet-MI & ba & 0.6377 $\pm$ 0.0042 & 0.5774 $\pm$ 0.0060 & +6.03 \\
PhysioNet-MI & kappa & 0.5169 $\pm$ 0.0056 & 0.4365 $\pm$ 0.0081 & +8.05 \\
PhysioNet-MI & f1 & 0.6386 $\pm$ 0.0041 & 0.5795 $\pm$ 0.0066 & +5.91 \\
FACED & ba & 0.4464 $\pm$ 0.0102 & 0.2842 $\pm$ 0.0109 & +16.22 \\
FACED & kappa & 0.3744 $\pm$ 0.0109 & 0.1961 $\pm$ 0.0119 & +17.83 \\
FACED & f1 & 0.4445 $\pm$ 0.0098 & 0.2865 $\pm$ 0.0106 & +15.80 \\
TUAB & ba & 0.8202 $\pm$ 0.0031 & 0.7759 $\pm$ 0.0057 & +4.43 \\
TUAB & pr-auc & 0.9091 $\pm$ 0.0034 & 0.8465 $\pm$ 0.0056 & +6.26 \\
TUAB & roc-auc & 0.9042 $\pm$ 0.0024 & 0.8514 $\pm$ 0.0051 & +5.29 \\
HMC & ba & 0.7476 $\pm$ 0.0032 & 0.6727 $\pm$ 0.0034 & +7.49 \\
HMC & kappa & 0.6987 $\pm$ 0.0033 & 0.6169 $\pm$ 0.0026 & +8.18 \\
HMC & f1 & 0.7634 $\pm$ 0.0017 & 0.6997 $\pm$ 0.0027 & +6.36 \\
TUEV & ba & 0.7261 $\pm$ 0.0135 & 0.4902 $\pm$ 0.0154 & +23.59 \\
TUEV & kappa & 0.7428 $\pm$ 0.0127 & 0.5070 $\pm$ 0.0219 & +23.58 \\
TUEV & f1 & 0.8655 $\pm$ 0.0055 & 0.7435 $\pm$ 0.0104 & +12.20 \\
ISRUC & ba & 0.8130 $\pm$ 0.0027 & 0.7675 $\pm$ 0.0036 & +4.55 \\
ISRUC & kappa & 0.7875 $\pm$ 0.0026 & 0.7373 $\pm$ 0.0041 & +5.03 \\
ISRUC & f1 & 0.8349 $\pm$ 0.0023 & 0.7955 $\pm$ 0.0034 & +3.94 \\
SHU-MI & ba & 0.6302 $\pm$ 0.0331 & 0.5864 $\pm$ 0.0115 & +4.38 \\
SHU-MI & pr-auc & 0.6891 $\pm$ 0.0457 & 0.6437 $\pm$ 0.0152 & +4.54 \\
SHU-MI & roc-auc & 0.6908 $\pm$ 0.0452 & 0.6171 $\pm$ 0.0185 & +7.37 \\
SEED-V & ba & 0.3930 $\pm$ 0.0035 & 0.2565 $\pm$ 0.0029 & +13.65 \\
SEED-V & kappa & 0.2442 $\pm$ 0.0045 & 0.0720 $\pm$ 0.0039 & +17.22 \\
SEED-V & f1 & 0.4016 $\pm$ 0.0038 & 0.2643 $\pm$ 0.0029 & +13.73 \\
BCIC2020-3 & ba & 0.6619 $\pm$ 0.0090 & 0.4725 $\pm$ 0.0184 & +18.93 \\
BCIC2020-3 & kappa & 0.5773 $\pm$ 0.0113 & 0.3407 $\pm$ 0.0230 & +23.67 \\
BCIC2020-3 & f1 & 0.6618 $\pm$ 0.0089 & 0.4727 $\pm$ 0.0184 & +18.91 \\
CHB-MIT & ba & 0.7204 $\pm$ 0.0175 & 0.6402 $\pm$ 0.0666 & +8.03 \\
CHB-MIT & pr-auc & 0.4759 $\pm$ 0.0118 & 0.3647 $\pm$ 0.0305 & +11.12 \\
CHB-MIT & roc-auc & 0.9116 $\pm$ 0.0071 & 0.8817 $\pm$ 0.0104 & +2.99 \\
MentalArithmetic & ba & 0.7507 $\pm$ 0.0464 & 0.6271 $\pm$ 0.0680 & +12.36 \\
MentalArithmetic & pr-auc & 0.7249 $\pm$ 0.0625 & 0.4189 $\pm$ 0.0777 & +30.60 \\
MentalArithmetic & roc-auc & 0.8498 $\pm$ 0.0433 & 0.6848 $\pm$ 0.0471 & +16.50 \\
\bottomrule\end{tabular}\end{table}

\paragraph{Discussion.}
Across the matched initialization tables (Tables~\ref{tab:app-init-effnet}--\ref{tab:app-init-vit}), visual pretraining is most beneficial
for TUEV, FACED, MentalArithmetic, and BCIC2020-3, while gains are smaller or
mixed on ISRUC, TUAB, and Mumtaz2016. The pattern also depends on the backbone: ConvNeXt-Tiny
and ViT-Small often show larger improvements than EfficientNet-B0, whereas some
tasks show little or negative change. These results support a transferable
visual initialization, but not a uniform replacement for EEG pretraining.

\subsection{Frozen linear-probe results}
\label{app:linear-probe}

For the frozen-transfer control, all backbone parameters are frozen and only
the dataset-specific linear prediction head is optimized. We evaluate
ConvNeXt-Tiny and ViT-Small across the same downstream datasets and five-seed
protocol used for the fine-tuning experiments. EfficientNet-B0 is excluded
from this control because its batch-normalization layers introduce an
additional choice about whether running statistics are updated or frozen;
this can confound a comparison intended to isolate representation quality.
The primary-metric results are summarized in Table~\ref{tab:linear-probe-results}; complete per-dataset metrics and seed aggregates are provided with the released experiment bundle.

\begin{table*}[t]\centering\scriptsize
\caption{Frozen linear-probe performance of DINOv3-initialized backbones. The backbone is frozen and only a linear task head is trained. Values are test means $\pm$ population standard deviations over five seeds.}
\label{tab:linear-probe-results}
\begin{tabular}{@{}lccc@{}}\toprule
Dataset & Metric & ConvNeXt-Tiny & ViT-Small \\ \midrule
CHB-MIT & pr\_auc & 0.4022$\pm$0.0035 & 0.4360$\pm$0.0042 \\
TUAB & pr\_auc & 0.8743$\pm$0.0008 & 0.8398$\pm$0.0003 \\
TUEV & kappa & 0.6653$\pm$0.0050 & 0.6887$\pm$0.0019 \\
ISRUC & kappa & 0.6087$\pm$0.0006 & 0.6043$\pm$0.0013 \\
HMC & kappa & 0.5698$\pm$0.0005 & 0.5810$\pm$0.0007 \\
FACED & kappa & 0.1863$\pm$0.0019 & 0.0996$\pm$0.0014 \\
SEED-V & kappa & 0.1027$\pm$0.0047 & 0.1297$\pm$0.0021 \\
PhysioNet-MI & kappa & 0.1969$\pm$0.0107 & 0.1320$\pm$0.0070 \\
SHU-MI & pr\_auc & 0.5136$\pm$0.0057 & 0.5012$\pm$0.0027 \\
BCIC2020-3 & kappa & 0.0303$\pm$0.0170 & 0.1373$\pm$0.0186 \\
Mumtaz2016 & pr\_auc & 0.9536$\pm$0.0014 & 0.9673$\pm$0.0004 \\
MentalArithmetic & pr\_auc & 0.6300$\pm$0.1552 & 0.5114$\pm$0.0690 \\
\bottomrule
\end{tabular}
\end{table*}

\begin{table}[tbp]
\centering\scriptsize
\setlength{\tabcolsep}{4pt}
\caption{Complete frozen linear-probe BA comparison. Reference columns follow the same fixed non-pooling choice as the main-text table.}
\label{tab:linear-probe-ba-context}
\begin{tabular}{@{}lcccc@{}}\toprule
Dataset & CBraMod$^\dagger$ & REVE-Base$^\dagger$ & ConvNeXt-Tiny & ViT-Small \\ \midrule
CHB-MIT & -- & -- & $0.7487\pm0.0042$ & $0.6809\pm0.0023$ \\
TUAB & $0.500\pm0.000$ & $0.809\pm0.004$ & $0.7908\pm0.0014$ & $0.7733\pm0.0007$ \\
TUEV & $0.482\pm0.037$ & $0.508\pm0.073$ & $0.5544\pm0.0077$ & $0.5806\pm0.0052$ \\
ISRUC & $0.430\pm0.043$ & $0.662\pm0.030$ & $0.6407\pm0.0006$ & $0.6397\pm0.0012$ \\
HMC & $0.538\pm0.009$ & $0.604\pm0.008$ & $0.6225\pm0.0006$ & $0.6357\pm0.0012$ \\
FACED & $0.261\pm0.013$ & $0.422\pm0.028$ & $0.2774\pm0.0017$ & $0.1953\pm0.0013$ \\
SEED-V & -- & -- & $0.2824\pm0.0036$ & $0.3020\pm0.0016$ \\
PhysioNet-MI & $0.531\pm0.015$ & $0.510\pm0.012$ & $0.3977\pm0.0080$ & $0.3491\pm0.0052$ \\
SHU-MI & -- & -- & $0.5180\pm0.0068$ & $0.5074\pm0.0045$ \\
BCIC2020-3 & $0.374\pm0.007$ & $0.390\pm0.017$ & $0.2243\pm0.0136$ & $0.3099\pm0.0149$ \\
Mumtaz2016 & $0.907\pm0.027$ & $0.931\pm0.021$ & $0.8702\pm0.0042$ & $0.8884\pm0.0007$ \\
MentalArithmetic & $0.605\pm0.020$ & $0.740\pm0.073$ & $0.7035\pm0.0986$ & $0.6208\pm0.0315$ \\
\bottomrule\end{tabular}
\tablenote{$^\dagger$Published non-pooling linear-probe columns from Table 4 of \citet{elouahidi2025reve}, including their CBraMod reproduction. Visual results are five-seed means $\pm$ population standard deviations. Shared benchmark partitions do not equalize readout or optimization protocols; these are contextual comparisons. Missing entries are unreported.}
\end{table}

\begin{table}[tbp]\centering\scriptsize
\setlength{\tabcolsep}{3pt}
\caption{Published linear-probe variants from REVE Table 4 on the nine overlapping datasets. B/L denotes REVE-Base/Large; P denotes pooling. Values and reported uncertainties are transcribed at their original precision \citep{elouahidi2025reve}. These references were not rerun locally.}
\label{tab:linear-probe-reference-variants}
\begin{tabular}{@{}lcccccc@{}}\toprule
Dataset & REVE-B (P) & REVE-B & REVE-L (P) & REVE-L & CBraMod (P) & CBraMod \\ \midrule
TUAB & $0.810\pm0.007$ & $0.809\pm0.004$ & $0.821\pm0.004$ & $0.809\pm0.004$ & $0.500\pm0.000$ & $0.500\pm0.000$ \\
TUEV & $0.592\pm0.008$ & $0.508\pm0.073$ & $0.630\pm0.003$ & $0.550\pm0.014$ & $0.219\pm0.009$ & $0.482\pm0.037$ \\
ISRUC & $0.697\pm0.011$ & $0.662\pm0.030$ & $0.743\pm0.004$ & $0.758\pm0.001$ & $0.407\pm0.049$ & $0.430\pm0.043$ \\
HMC & $0.647\pm0.008$ & $0.604\pm0.008$ & $0.703\pm0.003$ & $0.710\pm0.007$ & $0.368\pm0.001$ & $0.538\pm0.009$ \\
FACED & $0.240\pm0.010$ & $0.422\pm0.028$ & $0.283\pm0.003$ & $0.469\pm0.007$ & $0.117\pm0.005$ & $0.261\pm0.013$ \\
PhysioNet-MI & $0.537\pm0.005$ & $0.510\pm0.012$ & $0.551\pm0.001$ & $0.617\pm0.000$ & $0.256\pm0.002$ & $0.531\pm0.015$ \\
BCIC2020-3 & $0.234\pm0.009$ & $0.390\pm0.017$ & $0.274\pm0.001$ & $0.378\pm0.021$ & $0.214\pm0.003$ & $0.374\pm0.007$ \\
Mumtaz2016 & $0.962\pm0.003$ & $0.931\pm0.021$ & $0.985\pm0.006$ & $0.980\pm0.009$ & $0.859\pm0.009$ & $0.907\pm0.027$ \\
MentalArithmetic & $0.725\pm0.010$ & $0.740\pm0.073$ & $0.712\pm0.008$ & $0.665\pm0.103$ & $0.500\pm0.000$ & $0.605\pm0.020$ \\
\bottomrule\end{tabular}\end{table}

\subsection{Temporal folding versus contiguous reshape}\label{app:fold-reshape}
This control keeps the backbone, initialization, output dimensions, optimization
recipe, data split, checkpoint selection, and five-seed protocol fixed. Only the
mapping of samples to rows and columns changes: temporal folding uses
$X[c,wP+p]$, whereas the contiguous control groups consecutive chunks. Thus
both inputs contain exactly the same values, and $\Delta$ is temporal-folding
performance minus contiguous-reshape performance in percentage points. The full
results are grouped by backbone below.

\begin{table}[tbp]\centering\scriptsize\caption{Temporal folding versus contiguous reshape (efficientnet-b0).}\begin{tabular}{@{}l l c c c@{}}\toprule Dataset & Metric & Temporal fold & Chunk reshape & $\Delta$ \\ \midrule
ISRUC & ba & 0.8025 $\pm$ 0.0027 & 0.7691 $\pm$ 0.0055 & +3.35 \\
ISRUC & kappa & 0.7704 $\pm$ 0.0032 & 0.7397 $\pm$ 0.0044 & +3.08 \\
ISRUC & f1 & 0.8197 $\pm$ 0.0029 & 0.7938 $\pm$ 0.0040 & +2.59 \\
Mumtaz2016 & ba & 0.9153 $\pm$ 0.0061 & 0.9139 $\pm$ 0.0126 & +0.14 \\
Mumtaz2016 & pr-auc & 0.9864 $\pm$ 0.0045 & 0.9761 $\pm$ 0.0092 & +1.03 \\
Mumtaz2016 & roc-auc & 0.9846 $\pm$ 0.0053 & 0.9706 $\pm$ 0.0114 & +1.40 \\
FACED & ba & 0.5516 $\pm$ 0.0179 & 0.3958 $\pm$ 0.0579 & +15.58 \\
FACED & kappa & 0.4953 $\pm$ 0.0197 & 0.3198 $\pm$ 0.0649 & +17.55 \\
FACED & f1 & 0.5567 $\pm$ 0.0163 & 0.3999 $\pm$ 0.0604 & +15.68 \\
TUEV & ba & 0.6423 $\pm$ 0.0164 & 0.5826 $\pm$ 0.0315 & +5.97 \\
TUEV & kappa & 0.6901 $\pm$ 0.0190 & 0.6509 $\pm$ 0.0231 & +3.92 \\
TUEV & f1 & 0.8380 $\pm$ 0.0076 & 0.8176 $\pm$ 0.0085 & +2.04 \\
CHB-MIT & ba & 0.7552 $\pm$ 0.0599 & 0.7348 $\pm$ 0.0642 & +2.04 \\
CHB-MIT & pr-auc & 0.4092 $\pm$ 0.0536 & 0.3815 $\pm$ 0.0604 & +2.76 \\
CHB-MIT & roc-auc & 0.8974 $\pm$ 0.0194 & 0.8937 $\pm$ 0.0220 & +0.37 \\
MentalArithmetic & ba & 0.6910 $\pm$ 0.0418 & 0.6104 $\pm$ 0.0123 & +8.06 \\
MentalArithmetic & pr-auc & 0.7903 $\pm$ 0.0536 & 0.5849 $\pm$ 0.0734 & +20.54 \\
MentalArithmetic & roc-auc & 0.8779 $\pm$ 0.0251 & 0.7266 $\pm$ 0.0534 & +15.13 \\
TUAB & ba & 0.8317 $\pm$ 0.0032 & 0.8179 $\pm$ 0.0012 & +1.38 \\
TUAB & pr-auc & 0.9178 $\pm$ 0.0024 & 0.9031 $\pm$ 0.0016 & +1.47 \\
TUAB & roc-auc & 0.9126 $\pm$ 0.0024 & 0.8964 $\pm$ 0.0014 & +1.62 \\
SHU-MI & ba & 0.6273 $\pm$ 0.0085 & 0.6313 $\pm$ 0.0046 & -0.40 \\
SHU-MI & pr-auc & 0.6982 $\pm$ 0.0178 & 0.7009 $\pm$ 0.0138 & -0.28 \\
SHU-MI & roc-auc & 0.6962 $\pm$ 0.0139 & 0.6946 $\pm$ 0.0144 & +0.16 \\
HMC & ba & 0.7557 $\pm$ 0.0019 & 0.7056 $\pm$ 0.0039 & +5.01 \\
HMC & kappa & 0.7070 $\pm$ 0.0023 & 0.6379 $\pm$ 0.0046 & +6.91 \\
HMC & f1 & 0.7700 $\pm$ 0.0017 & 0.7168 $\pm$ 0.0041 & +5.32 \\
\bottomrule\end{tabular}\end{table}

\begin{table}[tbp]\centering\scriptsize\caption{Temporal folding versus contiguous reshape (convnext-tiny-dinov3).}\begin{tabular}{@{}l l c c c@{}}\toprule Dataset & Metric & Temporal fold & Chunk reshape & $\Delta$ \\ \midrule
ISRUC & ba & 0.8142 $\pm$ 0.0009 & 0.7894 $\pm$ 0.0016 & +2.48 \\
ISRUC & kappa & 0.7905 $\pm$ 0.0018 & 0.7650 $\pm$ 0.0014 & +2.56 \\
ISRUC & f1 & 0.8371 $\pm$ 0.0012 & 0.8162 $\pm$ 0.0012 & +2.10 \\
Mumtaz2016 & ba & 0.9176 $\pm$ 0.0138 & 0.9116 $\pm$ 0.0221 & +0.60 \\
Mumtaz2016 & pr-auc & 0.9842 $\pm$ 0.0020 & 0.9799 $\pm$ 0.0072 & +0.43 \\
Mumtaz2016 & roc-auc & 0.9825 $\pm$ 0.0034 & 0.9794 $\pm$ 0.0065 & +0.30 \\
FACED & ba & 0.6363 $\pm$ 0.0027 & 0.4079 $\pm$ 0.1514 & +22.85 \\
FACED & kappa & 0.5890 $\pm$ 0.0032 & 0.3321 $\pm$ 0.1694 & +25.69 \\
FACED & f1 & 0.6374 $\pm$ 0.0030 & 0.3926 $\pm$ 0.1809 & +24.47 \\
TUEV & ba & 0.6763 $\pm$ 0.0122 & 0.6312 $\pm$ 0.0524 & +4.51 \\
TUEV & kappa & 0.7095 $\pm$ 0.0305 & 0.6721 $\pm$ 0.0167 & +3.73 \\
TUEV & f1 & 0.8488 $\pm$ 0.0135 & 0.8259 $\pm$ 0.0121 & +2.29 \\
CHB-MIT & ba & 0.7681 $\pm$ 0.0294 & 0.7067 $\pm$ 0.0540 & +6.15 \\
CHB-MIT & pr-auc & 0.4734 $\pm$ 0.0345 & 0.4294 $\pm$ 0.0628 & +4.40 \\
CHB-MIT & roc-auc & 0.9067 $\pm$ 0.0156 & 0.9033 $\pm$ 0.0197 & +0.33 \\
MentalArithmetic & ba & 0.7667 $\pm$ 0.0234 & 0.7097 $\pm$ 0.0656 & +5.69 \\
MentalArithmetic & pr-auc & 0.7598 $\pm$ 0.0436 & 0.6198 $\pm$ 0.1237 & +14.01 \\
MentalArithmetic & roc-auc & 0.8752 $\pm$ 0.0313 & 0.7870 $\pm$ 0.0751 & +8.82 \\
TUAB & ba & 0.8216 $\pm$ 0.0050 & 0.7469 $\pm$ 0.1235 & +7.47 \\
TUAB & pr-auc & 0.9059 $\pm$ 0.0045 & 0.8603 $\pm$ 0.0649 & +4.56 \\
TUAB & roc-auc & 0.9020 $\pm$ 0.0045 & 0.8090 $\pm$ 0.1545 & +9.30 \\
SHU-MI & ba & 0.6358 $\pm$ 0.0075 & 0.6004 $\pm$ 0.0523 & +3.54 \\
SHU-MI & pr-auc & 0.6910 $\pm$ 0.0129 & 0.6822 $\pm$ 0.0333 & +0.88 \\
SHU-MI & roc-auc & 0.7006 $\pm$ 0.0112 & 0.6865 $\pm$ 0.0363 & +1.41 \\
HMC & ba & 0.7581 $\pm$ 0.0055 & 0.7244 $\pm$ 0.0039 & +3.37 \\
HMC & kappa & 0.7111 $\pm$ 0.0063 & 0.6700 $\pm$ 0.0050 & +4.11 \\
HMC & f1 & 0.7744 $\pm$ 0.0043 & 0.7406 $\pm$ 0.0036 & +3.38 \\
\bottomrule\end{tabular}\end{table}

\begin{table}[tbp]\centering\scriptsize\caption{Temporal folding versus contiguous reshape (vit-small-dinov3).}\begin{tabular}{@{}l l c c c@{}}\toprule Dataset & Metric & Temporal fold & Chunk reshape & $\Delta$ \\ \midrule
ISRUC & ba & 0.8130 $\pm$ 0.0027 & 0.7741 $\pm$ 0.0024 & +3.89 \\
ISRUC & kappa & 0.7875 $\pm$ 0.0026 & 0.7473 $\pm$ 0.0023 & +4.02 \\
ISRUC & f1 & 0.8349 $\pm$ 0.0023 & 0.8023 $\pm$ 0.0022 & +3.26 \\
Mumtaz2016 & ba & 0.9112 $\pm$ 0.0099 & 0.8857 $\pm$ 0.0320 & +2.55 \\
Mumtaz2016 & pr-auc & 0.9797 $\pm$ 0.0037 & 0.9729 $\pm$ 0.0094 & +0.68 \\
Mumtaz2016 & roc-auc & 0.9767 $\pm$ 0.0047 & 0.9642 $\pm$ 0.0187 & +1.25 \\
FACED & ba & 0.4464 $\pm$ 0.0102 & 0.4438 $\pm$ 0.0146 & +0.27 \\
FACED & kappa & 0.3744 $\pm$ 0.0109 & 0.3716 $\pm$ 0.0164 & +0.28 \\
FACED & f1 & 0.4445 $\pm$ 0.0098 & 0.4436 $\pm$ 0.0142 & +0.09 \\
TUEV & ba & 0.7261 $\pm$ 0.0135 & 0.6727 $\pm$ 0.0056 & +5.34 \\
TUEV & kappa & 0.7428 $\pm$ 0.0127 & 0.6970 $\pm$ 0.0074 & +4.58 \\
TUEV & f1 & 0.8655 $\pm$ 0.0055 & 0.8414 $\pm$ 0.0039 & +2.41 \\
CHB-MIT & ba & 0.7204 $\pm$ 0.0175 & 0.7614 $\pm$ 0.0118 & -4.09 \\
CHB-MIT & pr-auc & 0.4759 $\pm$ 0.0118 & 0.4958 $\pm$ 0.0123 & -2.00 \\
CHB-MIT & roc-auc & 0.9116 $\pm$ 0.0071 & 0.9262 $\pm$ 0.0056 & -1.46 \\
MentalArithmetic & ba & 0.7507 $\pm$ 0.0464 & 0.7118 $\pm$ 0.0655 & +3.89 \\
MentalArithmetic & pr-auc & 0.7249 $\pm$ 0.0625 & 0.7239 $\pm$ 0.0734 & +0.10 \\
MentalArithmetic & roc-auc & 0.8498 $\pm$ 0.0433 & 0.8435 $\pm$ 0.0381 & +0.63 \\
TUAB & ba & 0.8202 $\pm$ 0.0031 & 0.8136 $\pm$ 0.0022 & +0.66 \\
TUAB & pr-auc & 0.9091 $\pm$ 0.0034 & 0.8994 $\pm$ 0.0017 & +0.96 \\
TUAB & roc-auc & 0.9042 $\pm$ 0.0024 & 0.8955 $\pm$ 0.0010 & +0.87 \\
SHU-MI & ba & 0.6302 $\pm$ 0.0331 & 0.5519 $\pm$ 0.0652 & +7.83 \\
SHU-MI & pr-auc & 0.6891 $\pm$ 0.0457 & 0.5996 $\pm$ 0.1003 & +8.94 \\
SHU-MI & roc-auc & 0.6908 $\pm$ 0.0452 & 0.5972 $\pm$ 0.0910 & +9.36 \\
HMC & ba & 0.7476 $\pm$ 0.0032 & 0.6751 $\pm$ 0.0060 & +7.26 \\
HMC & kappa & 0.6987 $\pm$ 0.0033 & 0.6139 $\pm$ 0.0053 & +8.48 \\
HMC & f1 & 0.7634 $\pm$ 0.0017 & 0.6974 $\pm$ 0.0048 & +6.60 \\
\bottomrule\end{tabular}\end{table}

\paragraph{Discussion.}
The temporal-folding control shows a consistent advantage for CNN backbones across
most datasets, with the largest differences on FACED and MentalArithmetic and
smaller but positive differences on clinical and sleep tasks (the only CNN
exception is a 0.3pp SHU-MI reversal for EfficientNet-B0). ViT-Small is more
task-sensitive: it benefits clearly on TUEV, HMC, and SHU-MI, but contiguous
reshape is better on CHB-MIT and nearly tied on FACED and MentalArithmetic.
Thus the gain is not explained by height expansion alone; it depends on how the
backbone uses the phase-structured neighborhoods.
\subsection{Fold-factor sensitivity}
\label{app:p-ablation}

To examine whether the preferred fold factor is task- and backbone-dependent,
we run controlled sweeps with $P\in\{1,2,4,8\}$ on CHB-MIT, TUEV, and
MentalArithmetic. All other training settings, data splits, validation
checkpoint selection, and final test protocols are held fixed within each
backbone; the $P=4$ rows coincide with the finalized recipes used in the main
tables. Bold and underlined values are the best and second-best values within
each backbone and metric, respectively.

\begin{table*}[t]
\centering
\caption{CHB-MIT fold-factor sensitivity (BA / PR-AUC / ROC-AUC). Values are
test means $\pm$ population standard deviations after validation-based
checkpoint selection.}
\label{tab:p-ablation-chb}
\scriptsize
\setlength{\tabcolsep}{3pt}
\renewcommand{\arraystretch}{1.12}
\begin{tabular*}{\textwidth}{@{\extracolsep{\fill}}lccccc@{}}
\toprule
Backbone & $P$ & Folded shape & BA & PR-AUC & ROC-AUC \\
\midrule
EfficientNet-B0 & 1 & $16\times2000$ & $0.7136\pm0.0425$ & $0.3686\pm0.0958$ & \bestcell{0.9171\pm0.0089} \\
 & 2 & $32\times1000$ & $0.6913\pm0.0357$ & $0.3676\pm0.0498$ & $0.8772\pm0.0409$ \\
 & 4 & $64\times500$ & \secondcell{0.7552\pm0.0599} & \bestcell{0.4092\pm0.0536} & \secondcell{0.8974\pm0.0194} \\
 & 8 & $128\times250$ & \bestcell{0.8210\pm0.0249} & \secondcell{0.3879\pm0.0493} & $0.8941\pm0.0133$ \\
\midrule
ConvNeXt-Tiny & 1 & $16\times2000$ & $0.7170\pm0.0100$ & \bestcell{0.5103\pm0.0057} & $0.9078\pm0.0046$ \\
 & 2 & $32\times1000$ & $0.7249\pm0.0306$ & \secondcell{0.4960\pm0.0106} & \bestcell{0.9112\pm0.0054} \\
 & 4 & $64\times500$ & \bestcell{0.7681\pm0.0294} & $0.4734\pm0.0345$ & $0.9066\pm0.0155$ \\
 & 8 & $128\times250$ & \secondcell{0.7391\pm0.0430} & $0.4264\pm0.0334$ & \secondcell{0.9084\pm0.0118} \\
\midrule
ViT-Small & 1 & $16\times2000$ & $0.6884\pm0.0115$ & $0.3527\pm0.0482$ & \secondcell{0.9101\pm0.0067} \\
 & 2 & $32\times1000$ & $0.6734\pm0.0322$ & $0.4235\pm0.0303$ & $0.9065\pm0.0084$ \\
 & 4 & $64\times500$ & \secondcell{0.7204\pm0.0175} & \bestcell{0.4759\pm0.0118} & \bestcell{0.9116\pm0.0071} \\
 & 8 & $128\times250$ & \bestcell{0.7540\pm0.0224} & \secondcell{0.4510\pm0.0159} & $0.9094\pm0.0078$ \\
\bottomrule
\end{tabular*}
\tablenote{Each row is a separate five-seed run; the $P=4$ row matches the
finalized CHB-MIT recipe. The primary metric is PR-AUC.}
\end{table*}

\begin{table*}[t]
\centering
\caption{TUEV fold-factor sensitivity (BA / Cohen's $\kappa$ / weighted F1).
Values are test means $\pm$ population standard deviations after
validation-based checkpoint selection.}
\label{tab:p-ablation-tuev}
\scriptsize
\setlength{\tabcolsep}{3pt}
\renewcommand{\arraystretch}{1.12}
\begin{tabular*}{\textwidth}{@{\extracolsep{\fill}}lccccc@{}}
\toprule
Backbone & $P$ & Folded shape & BA & $\kappa$ & Weighted F1 \\
\midrule
EfficientNet-B0 & 1 & $16\times1000$ & $0.5172\pm0.0679$ & $0.6298\pm0.0235$ & $0.8034\pm0.0161$ \\
 & 2 & $32\times500$ & \secondcell{0.5859\pm0.0871} & \bestcell{0.6907\pm0.0473} & \secondcell{0.8355\pm0.0253} \\
 & 4 & $64\times250$ & \bestcell{0.6423\pm0.0164} & \secondcell{0.6901\pm0.0190} & \bestcell{0.8380\pm0.0076} \\
 & 8 & $128\times125$ & $0.5583\pm0.0878$ & $0.6308\pm0.0429$ & $0.8118\pm0.0215$ \\
\midrule
ConvNeXt-Tiny & 1 & $16\times1000$ & $0.4977\pm0.1376$ & $0.5381\pm0.1921$ & $0.7576\pm0.0965$ \\
 & 2 & $32\times500$ & $0.6274\pm0.0263$ & $0.7042\pm0.0414$ & $0.8475\pm0.0197$ \\
 & 4 & $64\times250$ & \secondcell{0.6763\pm0.0122} & \secondcell{0.7095\pm0.0305} & \secondcell{0.8488\pm0.0135} \\
 & 8 & $128\times125$ & \bestcell{0.7001\pm0.0232} & \bestcell{0.7274\pm0.0124} & \bestcell{0.8576\pm0.0084} \\
\midrule
ViT-Small & 1 & $16\times1000$ & $0.6577\pm0.0108$ & $0.6848\pm0.0063$ & $0.8311\pm0.0029$ \\
 & 2 & $32\times500$ & $0.6962\pm0.0153$ & \secondcell{0.7150\pm0.0168} & \secondcell{0.8514\pm0.0098} \\
 & 4 & $64\times250$ & \secondcell{0.7261\pm0.0135} & \bestcell{0.7428\pm0.0127} & \bestcell{0.8655\pm0.0055} \\
 & 8 & $128\times125$ & \bestcell{0.7383\pm0.0077} & $0.7353\pm0.0079$ & $0.8638\pm0.0037$ \\
\bottomrule
\end{tabular*}
\tablenote{Each row is a separate five-seed run; the $P=4$ row matches the
finalized TUEV recipe. The primary metric is $\kappa$.}
\end{table*}

\begin{table*}[t]
\centering
\caption{MentalArithmetic fold-factor sensitivity (BA / PR-AUC / ROC-AUC).
Values are test means $\pm$ population standard deviations after
validation-based checkpoint selection.}
\label{tab:p-ablation-mental}
\scriptsize
\setlength{\tabcolsep}{3pt}
\renewcommand{\arraystretch}{1.12}
\begin{tabular*}{\textwidth}{@{\extracolsep{\fill}}lccccc@{}}
\toprule
Backbone & $P$ & Folded shape & BA & PR-AUC & ROC-AUC \\
\midrule
EfficientNet-B0 & 1 & $20\times1000$ & $0.4764\pm0.0499$ & $0.2119\pm0.0839$ & $0.3150\pm0.0931$ \\
 & 2 & $40\times500$ & \bestcell{0.7056\pm0.0207} & \secondcell{0.6672\pm0.0409} & \secondcell{0.8012\pm0.0365} \\
 & 4 & $80\times250$ & \secondcell{0.6910\pm0.0418} & \bestcell{0.7903\pm0.0536} & \bestcell{0.8779\pm0.0251} \\
 & 8 & $160\times125$ & $0.6646\pm0.0341$ & $0.5944\pm0.0578$ & $0.7411\pm0.0332$ \\
\midrule
ConvNeXt-Tiny & 1 & $20\times1000$ & $0.6215\pm0.1549$ & $0.6968\pm0.1400$ & $0.8604\pm0.0920$ \\
 & 2 & $40\times500$ & \secondcell{0.7514\pm0.0970} & \bestcell{0.7781\pm0.1431} & \bestcell{0.8991\pm0.0905} \\
 & 4 & $80\times250$ & \bestcell{0.7667\pm0.0234} & \secondcell{0.7598\pm0.0436} & \secondcell{0.8752\pm0.0313} \\
 & 8 & $160\times125$ & $0.6326\pm0.0284$ & $0.4703\pm0.0694$ & $0.6825\pm0.0351$ \\
\midrule
ViT-Small & 1 & $20\times1000$ & \secondcell{0.7618\pm0.0282} & \secondcell{0.7814\pm0.0259} & \secondcell{0.8942\pm0.0143} \\
 & 2 & $40\times500$ & \bestcell{0.7792\pm0.0442} & \bestcell{0.7975\pm0.0383} & \bestcell{0.8956\pm0.0306} \\
 & 4 & $80\times250$ & $0.7507\pm0.0464$ & $0.7249\pm0.0625$ & $0.8498\pm0.0433$ \\
 & 8 & $160\times125$ & $0.6500\pm0.0183$ & $0.5565\pm0.0394$ & $0.7526\pm0.0223$ \\
\bottomrule
\end{tabular*}
\tablenote{Each row is a separate five-seed run; the $P=4$ row matches the
finalized MentalArithmetic recipe. The primary metric is PR-AUC.}
\end{table*}

The sweeps confirm that the preferred factor is task- and backbone-dependent.
On MentalArithmetic, EfficientNet-B0 peaks at $P=4$ whereas both DINOv3
backbones peak at $P=2$ on PR-AUC. On CHB-MIT, PR-AUC is best at $P=4$ for
EfficientNet-B0 and ViT-Small but at $P=1$ for ConvNeXt-Tiny, and
EfficientNet-B0 reaches its highest BA at $P=8$ without improving its
validation-selected metric. On TUEV, $\kappa$ and weighted F1 are best at
$P=8$ for ConvNeXt-Tiny and at $P=4$ for ViT-Small, while EfficientNet-B0 is
nearly tied between $P=2$ and $P=4$ ($\kappa=.6907$ vs $.6901$). The geometry
rule therefore serves as a reproducible default; validation-based tuning can
improve the result when additional compute is available.


\section{Model-Family Comparison}
\label{app:model-families}

This appendix separates the representational interface from model scale. All
three vision backbones receive the same folded EEG tensor, whereas the EEG
foundation models use their native patching and positional mechanisms. The
comparison is descriptive rather than a runtime benchmark: parameter and FLOP
counts are measured on a shared input and exclude task heads. Using the
released implementations and an explicit PyTorch FLOP counter, the vision
encoders at a shared $32\times1000$ input with $P=4$ require 4.01M/0.247G
(EfficientNet-B0), 27.82M/2.903G (ConvNeXt-Tiny), and 21.39M/3.000G
(ViT-Small) parameters/FLOPs, respectively. FLOPs count one multiply-add as
one operation; attention fast paths are disabled for reproducible accounting.

\begin{table*}[tbp]
\centering
\caption{Conceptual scope of the model families discussed in the paper.}
\label{tab:model-scale}
\small
\renewcommand{\arraystretch}{1.06}
\begin{tabularx}{\textwidth}{@{}lXXX@{}}
\toprule
 & \method{} (B0) & \cbramod{} & \reve{}-Base \\
\midrule
Input & phase-folded raw waveform & raw patches + FFT magnitude & raw waveform patches \\
Encoder & EfficientNet-B0 & criss-cross transformer & coordinate-aware transformer \\
EEG pretraining & none & TUH EEG & 92 EEG datasets \\
Encoder parameters & 4.01M & 4.92M & 69.19M \\
Forward compute @$32\times1000$ & 0.25 GFLOPs & 0.79 GFLOPs & 11.07 GFLOPs \\
Spatial identity & row order & learned fixed-order encoding & physical coordinates \\
Main target & full fine-tuning & compact EEG transfer & montage-flexible transfer \\
\bottomrule
\end{tabularx}
\tablenote{The B0 encoder count excludes its unused ImageNet classifier and
the task-specific head. Parameter and compute counts for \cbramod{} and
\reve{}-Base are measured from their released implementations at the shared
$32\times1000$ input; the \cbramod{} paper reports about 4.0M parameters,
whereas its released checkpoint contains 4.92M.}
\end{table*}

\section{Training Recipes and Backbone Input Shapes}
\label{app:recipes}

Table~\ref{tab:recipes} summarizes one representative ConvNeXt-Tiny training
recipe for each dataset. ``Folded'' denotes the non-padded image shape; zero
padding is appended only at the bottom and right. All runs use DINOv3
initialization, AdamW, cosine decay, dropout .1, bfloat16 mixed precision, and
a dataset-specific linear head. Learning rates, epoch budgets, fold factors,
and validation metrics are selected before the final test evaluation.

\begin{table*}[tbp]
\centering
\caption{Finalized phase-interleaved downstream recipes.}
\label{tab:recipes}
\footnotesize
\setlength{\tabcolsep}{4pt}
\renewcommand{\arraystretch}{1.06}
\begin{tabularx}{\textwidth}{@{}lccrccc@{}}
\toprule
Dataset & $P$ & Folded $H\times W$ & Epochs & LR & WD & Select \\
\midrule
CHB-MIT & 4 & $64\times500$ & 30 & $10^{-4}$ & $5\cdot10^{-4}$ & PR-AUC \\
TUAB & 4 & $64\times500$ & 10 & $10^{-4}$ & $5\cdot10^{-4}$ & PR-AUC \\
TUEV & 4 & $64\times250$ & 5 & $10^{-4}$ & $5\cdot10^{-4}$ & $\kappa$ \\
ISRUC & 12 & $72\times500$ & 15 & $10^{-4}$ & $5\cdot10^{-3}$ & $\kappa$ \\
FACED & 2 & $64\times1000$ & 50 & $10^{-4}$ & $5\cdot10^{-4}$ & $\kappa$ \\
SEED-V & 1 & $62\times200$ & 50 & $10^{-4}$ & $5\cdot10^{-3}$ & $\kappa$ \\
PhysioNet-MI & 1 & $64\times800$ & 30 & $10^{-4}$ & $5\cdot10^{-3}$ & $\kappa$ \\
SHU-MI & 2 & $64\times400$ & 10 & $10^{-4}$ & $5\cdot10^{-4}$ & PR-AUC \\
BCIC2020-3 & 1 & $64\times600$ & 70 & $10^{-5}$ & $5\cdot10^{-3}$ & $\kappa$ \\
Mumtaz2016 & 4 & $76\times250$ & 30 & $10^{-4}$ & $5\cdot10^{-4}$ & PR-AUC \\
MentalArithmetic & 4 & $80\times250$ & 50 & $10^{-4}$ & $5\cdot10^{-4}$ & PR-AUC \\
HMC & 16 & $64\times375$ & 30 & $10^{-4}$ & $5\cdot10^{-3}$ & $\kappa$ \\
\bottomrule
\end{tabularx}
\tablenote{``Select'' is the validation metric used for checkpointing.
ConvNeXt-Tiny uses flatten aggregation except for TUEV, ISRUC, and HMC, which
use GAP.}
\end{table*}

\section{Data and Code Reproducibility Statement}

The codebase records exact dataset paths, input shapes, preprocessing, seeds,
training arguments, logs, and validation-selected checkpoints.  Public data
licenses and access conditions differ by benchmark; no dataset is redistributed
with the paper.  The CHB-MIT manifest follows the subject split of the
\cbramod{} benchmark code (CHB01--20/CHB21--22/CHB23--24 for
train/validation/test) \citep{wang2024cbramod}; the remaining manifests follow
the fixed subject- or trial-level partitions of the \cbramod{} and \reve{}
benchmark code \citep{wang2024cbramod,elouahidi2025reve}, with HMC using the
\reve{} partition \citep{elouahidi2025reve,alvarez2022hmc}.  Reproduction still
requires the documented dataset versions and local split manifests. We use
TUAB v3.0.1 and TUEV v2.0.1; for TUEV, the held-out test partition is the
official evaluation set and contains 29,421 segments in both v2.0.0 and
v2.0.1, while the newer release changes only the training and validation
counts (68,445/15,487 rather than 67,436/15,634).
